\documentclass[11pt]{article}
\usepackage[a4paper,margin=28mm,headheight=14pt]{geometry}
\usepackage[T1]{fontenc}
\usepackage[utf8]{inputenc}
\usepackage{lmodern,microtype}
\usepackage{amsmath,amssymb,amsthm,mathtools}
\usepackage{booktabs,array}
\usepackage{enumitem}
\usepackage{fancyhdr}
\usepackage{aliascnt,placeins}
\usepackage[hidelinks]{hyperref}
\usepackage[nameinlink,noabbrev]{cleveref}
\hypersetup{pdftitle={Exact Orders, Shift Optimality, and Zero Asymptotics for q-Exponential Hankel Determinants},pdfauthor={KEIO},pdfsubject={Exact orders, arithmetic shift bounds, and analytic zero asymptotics}}
\pagestyle{fancy}
\fancyhf{}
\fancyhead[L]{\small KEIO}
\fancyhead[R]{\small Shifted $q$-exponential Hankel determinants}
\fancyfoot[C]{\thepage}
\renewcommand{\headrulewidth}{0.3pt}
\setlength{\emergencystretch}{2em}
\setlist[enumerate]{leftmargin=*,itemsep=3pt,topsep=5pt}
\numberwithin{equation}{section}
\newtheorem{theorem}{Theorem}[section]
\newaliascnt{proposition}{theorem}
\newtheorem{proposition}[proposition]{Proposition}
\aliascntresetthe{proposition}
\newaliascnt{lemma}{theorem}
\newtheorem{lemma}[lemma]{Lemma}
\aliascntresetthe{lemma}
\newaliascnt{corollary}{theorem}
\newtheorem{corollary}[corollary]{Corollary}
\aliascntresetthe{corollary}
\theoremstyle{definition}
\newaliascnt{definition}{theorem}
\newtheorem{definition}[definition]{Definition}
\aliascntresetthe{definition}
\theoremstyle{remark}
\newaliascnt{remark}{theorem}
\newtheorem{remark}[remark]{Remark}
\aliascntresetthe{remark}
\crefname{proposition}{Proposition}{Propositions}
\Crefname{proposition}{Proposition}{Propositions}
\crefname{lemma}{Lemma}{Lemmas}
\Crefname{lemma}{Lemma}{Lemmas}
\crefname{corollary}{Corollary}{Corollaries}
\Crefname{corollary}{Corollary}{Corollaries}
\crefname{remark}{Remark}{Remarks}
\Crefname{remark}{Remark}{Remarks}
\crefname{definition}{Definition}{Definitions}
\Crefname{definition}{Definition}{Definitions}
\newcommand{\Q}{\mathbb{Q}}
\newcommand{\Z}{\mathbb{Z}}
\newcommand{\C}{\mathbb{C}}
\newcommand{\Qbar}{\overline{\mathbb{Q}}}
\newcommand{\ordq}{\operatorname{ord}_q}
\newcommand{\degq}{\operatorname{deg}_q}
\newcommand{\Vnr}{V_n^{(r)}}
\newcommand{\Fq}{E_q}
\newcommand{\DD}{\mathcal{D}}
\newcommand{\hplus}[1]{\left(#1\right)_+}
\newcommand{\fl}[1]{\left\lfloor#1\right\rfloor}
\newcommand{\prodempty}{\text{empty products are }1}
\title{\bfseries Exact Orders, Shift Optimality, and Zero Asymptotics\\[2mm]
for $q$-Exponential Hankel Determinants}
\author{KEIO}
\date{Revised manuscript, version 2.0\\11 October 2026}
\begin{document}
\maketitle
\begin{abstract}
We study shifted Hankel determinants of the remainder sequence
$v_0=\mu-1$, $v_k=(q^k-1)v_{k-1}-\alpha^k$ for the $q$-exponential.
A two-node formal moment expansion gives the exact generic order at $q=0$,
explicit lowest coefficients, and the polynomial degree.
At $\mu=E_q(\alpha)$, positive geometric moments give sharp size estimates.
A degree budget for a prescribed cyclotomic divisor proves that proportional
index shifts cannot improve the critical values $12/7$ and $6$ in the
associated cubic norm comparison; these arithmetic values are already known.
For real $q>1$, we then extend the moment representation to complex
$\alpha$, remove the exact forced zeros, and obtain a Schur expansion
of the normalized entire determinant. After scaling by $q^{n+r}$, this
function converges locally uniformly, uniformly in all shifts $r\ge0$, to
$\prod_{\ell\ge1}(1+zq^{-\ell})^\ell$, with an explicit geometric error.
This gives a full leading asymptotic and its first correction at fixed
$\alpha$, and determines the limiting zero clusters.
Using classical $q$-Meixner polynomials, we reduce a special shifted
determinant from size $n$ to size $r$ and locate the first additional zero
to relative error $O(q^{-n})$ for every fixed shift. We also resolve the
second unshifted cluster into two real simple zeros with distinct
superexponential displacement scales.
These analytic refinements do not assert a lower general irrationality
threshold. Modular specializations are recorded separately as consequences
of existing transcendence theory.
\end{abstract}
\noindent\textbf{Keywords.} $q$-exponential; Hankel determinant; cyclotomic divisibility;
Schur functions; $q$-Meixner polynomials; zero asymptotics.

\section{Introduction and principal results}\label{sec:introduction}

For $|q|>1$, consider the entire function
\begin{equation}\label{eq:exponential}
 E_q(z)=F_q(z;1)=\sum_{k=0}^{\infty}
 \frac{z^k}{\prod_{j=1}^{k}(q^j-1)}.
\end{equation}
Krattenthaler, Rochev, V\"a\"an\"anen, and Zudilin
\cite{KRVZ} use Hankel determinants of remainder sequences to study the
arithmetic of this function and related $q$-series. Their construction separates
three cubic growth quantities: smallness at the distinguished value, degree of
a common polynomial divisor, and degree or conjugate growth of the determinant.
For the remainder sequence used here, these quantities are especially amenable
to an exact analysis of an index shift.

Throughout the algebraic sections, $q,\alpha,\mu$ are independent
indeterminates, and
\begin{equation}\label{eq:recurrence}
 R=\Z[q,\alpha,\mu],\qquad
 v_0=\mu-1,\qquad v_k=(q^k-1)v_{k-1}-\alpha^k\quad(k\ge1).
\end{equation}
For integers $n\ge1$ and $r\ge0$, put
\begin{equation}\label{eq:hankel}
 V_n^{(r)}=\det(v_{i+j+r})_{0\le i,j<n},\qquad V_n=V_n^{(0)}.
\end{equation}
At the specialization $\mu=E_q(\alpha)$, direct solution of
\eqref{eq:recurrence} gives
\begin{equation}\label{eq:tail}
 v_k=\sum_{\ell=k+1}^{\infty}
 \frac{\alpha^\ell}{\prod_{j=k+1}^{\ell}(q^j-1)}.
\end{equation}
We use \eqref{eq:tail} analytically only where its parameters are specified.
In particular, the positivity and exact smallness results below assume
$q>1$ and $\alpha>0$.

The first part of this article gives complete proofs of the exact algebraic
endpoints, the prescribed cyclotomic divisor, and the restriction on the
associated shift comparison. The second part studies the determinant as an
entire function of $\alpha$ and resolves zero displacements that the cubic
arithmetic estimates cannot detect. The two parts address different scales:
the arithmetic comparison uses terms of order $n^3$, whereas the analytic
refinements determine constants, geometric corrections, and exponentially
small separations of zeros.

The thresholds $12/7$ and $6$ are not new arithmetic bounds. Rochev
\cite[Corollary 1.1]{Rochev2011} already obtains the corresponding
$7/12$ and $1/6$ conditions in his height notation. Our optimality statement
is restricted to the divisor and norm comparison specified below.
The analytic arguments use classical Schur identities \cite{Macdonald}
and $q$-Meixner orthogonality \cite{KS}; these inputs are explicitly
distinguished from the determinant identities and zero estimates proved here.

\begin{theorem}[Exact algebraic endpoints]\label{thm:main-algebra}
For all $n\ge1$ and $r\ge0$, in $\Q(\alpha,\mu)[q]$ one has
\begin{align}
 e_{n,r}:=\ordq V_n^{(r)}
 &=\binom n3+r\fl{\frac{(n-1)^2}{4}},\label{eq:exact-order}\\
 D_{n,r}:=\degq V_n^{(r)}
 &=\sum_{i=0}^{n-1}\frac{(2i+r)(2i+r+1)}2.\label{eq:exact-degree}
\end{align}
The lowest coefficients are given explicitly in \cref{prop:leading}.
Consequently, $q^{e_{n,r}}$ divides $V_n^{(r)}$ in $R$.
The equality in \eqref{eq:exact-order} is generic: after specializing
$\alpha$ or $\mu$, the order may increase.
\end{theorem}

\begin{theorem}[Positive moments and exact size]\label{thm:main-positive}
Let $q>1$, $\alpha>0$, and $\mu=E_q(\alpha)$. Define
\begin{equation}\label{eq:positive-nodes}
 x_j=\alpha q^{-j},\qquad
 w_j=\alpha q^{-j}\prod_{\ell=1}^{j-1}(1+\alpha q^{-\ell})
 \quad(j\ge1).
\end{equation}
Then $v_k=\sum_{j\ge1}w_jx_j^k$ for every $k\ge0$, and
$V_n^{(r)}>0$ for every $n\ge1,r\ge0$. With
\begin{equation}\label{eq:KZ}
 K_{n,r}=\frac{n(n+1)(2n+1)}6+\frac{rn(n+1)}2,
 \qquad Z_{n,r}=\alpha^{n^2+rn}q^{-K_{n,r}},
\end{equation}
we have
\begin{equation}\label{eq:main-bounds}
 Z_{n,r}\prod_{k=1}^{n-1}(1-q^{-k})^{2(n-k)}
 \le V_n^{(r)}
 \le Z_{n,r}\mu^n\prod_{m=1}^n(1-q^{-m(m+r)})^{-1}.
\end{equation}
In particular, uniformly in integers $r\ge0$,
\begin{equation}\label{eq:log-asymptotic}
 \log V_n^{(r)}=-K_{n,r}\log q+(n^2+rn)\log\alpha+O_{q,\alpha}(n).
\end{equation}
\end{theorem}

To state the third result, we specify the divisor rather than invoking an
unspecified ``full content.'' Write
\begin{align}
 S_d(n)&=\sum_{i=0}^{n-1}\fl{\frac{i}{2d}},\label{eq:S}\\
 T_d(n,p)&=\sum_{i=0}^{n-1}
 \min\left\{\fl{\frac{i}{d}},\fl{\frac{i+p}{2d}}\right\},\label{eq:T}\\
 \nu_d(n,r)&=\max_{p+p'=r}\bigl(T_d(n,p)+T_d(n,p')\bigr),\label{eq:nu}\\
 \DD_{n,r}(q)&=q^{e_{n,r}}\prod_{d\ge1}\Phi_d(q)^{\nu_d(n,r)},\label{eq:divisor}
\end{align}
where $p,p'$ are nonnegative integers and $\Phi_d$ is the $d$-th
cyclotomic polynomial. The product is finite because $\nu_d(n,r)=0$
for $d\ge n$.

\begin{theorem}[Optimality in the prescribed shift comparison]\label{thm:main-optimality}
The polynomial $\DD_{n,r}$ divides $V_n^{(r)}$ in $R$. Fix $t\ge0$ and
$r=\lfloor tn\rfloor$. The exact smallness and degree constants are
\begin{equation}\label{eq:AC}
 A(t)=\frac13+\frac t2,\qquad C(t)=\frac23+t+\frac{t^2}{2}.
\end{equation}
Every accumulation point $B$ of $n^{-3}\degq\DD_{n,\lfloor tn\rfloor}$
satisfies
\begin{equation}\label{eq:Bbar}
 B\le\overline B(t):=\frac14+\frac{3t}{8}+\frac{t^2}{8}.
\end{equation}
For such a $B$, the rationality-exclusion critical value of the standard
cubic comparison satisfies
\begin{equation}\label{eq:tau-bound}
 \tau_1(t;B):=\frac{A(t)+C(t)}{A(t)+B}
 \ge\frac{24+36t+12t^2}{14+21t+3t^2}\ge\frac{12}{7}.
\end{equation}
The last inequality is strict for $t>0$. For quadratic-field exclusion,
\begin{equation}\label{eq:tau2-bound}
 \tau_2(t;B):=\frac{A(t)+C(t)}{A(t)-C(t)+2B}\ge6
\end{equation}
whenever its denominator is positive. At $t=0$, $B=1/4$, and both
unshifted critical values are attained.
\end{theorem}

The term \emph{critical value} in this theorem refers to a particular
sufficient-condition calculation, defined by \eqref{eq:tau-bound} and
\eqref{eq:tau2-bound}. It is not a boundary between rational and irrational
values of $E_q(\alpha)$. Extra arithmetic divisors, specializations, alternative
determinants, and stronger conjugate estimates are outside this comparison.
Using accumulation points in the statement avoids assuming the existence of a
shifted divisor-degree limit that is not needed for the argument.

\subsection{Analytic refinements and their scope}

For the remainder of this introduction let $q>1$, $p=q^{-1}$, and write
$H_{n,r}(\alpha)=V_n^{(r)}(q,\alpha,E_q(\alpha))$.
After removal of its forced factors
\[
 \mathcal L_{n,r}(\alpha)
 =\alpha^{n^2+rn}\prod_{m=1}^{n-1}(1+\alpha p^m)^{n-m},
\]
we normalize by a positive geometric Cauchy determinant
$\mathcal C_{n,r}$; see \eqref{eq:analytic-normalization}.
The quotient $\Psi_{n,r}=H_{n,r}/(\mathcal L_{n,r}\mathcal C_{n,r})$
is entire and equals one at the origin. The following are the principal
analytic conclusions of this revision.

\begin{enumerate}
\item \textbf{A uniform entire scaling limit.}
For every $R>0$,
\[
 \sup_{r\ge0,\,|z|\le R}
 \left|\Psi_{n,r}(p^{-n-r}z)-\Phi_p(z)\right|=O_{p,R}(p^n),
 \qquad
 \Phi_p(z)=\prod_{\ell\ge1}(1+zp^\ell)^\ell.
\]
Thus the $k$th limiting cluster has exactly $k$ zeros counted with
multiplicity. At fixed $\alpha$, the same estimate yields an explicit
leading product and its first correction
(\cref{thm:scaled-limit,thm:fixed-alpha-refinement}).

\item \textbf{The first additional zero for every fixed shift.}
Let $\pi_r=\prod_{j=1}^r(1-p^j)$ and
$P_\infty=\prod_{j\ge1}(1-p^j)$. For each fixed $r\ge0$, the nearest
zero of $\Psi_{n,r}$ is real, simple, and satisfies
\[
 \beta_{n,r}
 =-q^{n+r+1}
 -\frac{p^{n^2-n-r-1}}{P_\infty\pi_r}
   \left(1+O_{p,r}(p^n)\right).
\]
An exact $r\times r$ determinant identity supplies the displacement
constant (\cref{sz:reduction,sz:first-zero}).

\item \textbf{Asymmetric splitting of the second unshifted cluster.}
For all sufficiently large $n$, the two zeros near $-q^{n+2}$ are
real and simple. Writing them as
$-q^{n+2}-\delta_n^-$ and $-q^{n+2}+\delta_n^+$, one has
\[
 \delta_n^-\sim
 \frac{p^{n^2-3n-1}}{(1-p)^2P_\infty},
 \qquad
 \delta_n^+\sim
 \frac{(1-p)p^{n^2+n-2}}{P_\infty}.
\]
Their separation scales follow from an exact product at $\mu=0$,
a trace identity, and classical $q$-Meixner polynomials
(\cref{mx:double-splitting}).
\end{enumerate}

All asymptotic statements hold for fixed $q>1$. Uniformity in $r$ is
asserted for the scaling limit, whereas the precise displacement of the
first zero holds for each fixed $r$. We do not infer reality or individual
splitting laws for all higher clusters from their limiting multiplicities.
These conclusions refine the analytic structure of the same determinants;
they do not change the cubic arithmetic exponents or prove a new general
irrationality criterion.

\section{Formal moments, exact orders, and extreme coefficients}\label{sec:formal}

\subsection{A two-node formal moment expansion}

Set $K=\Q(\alpha,\mu)$ and work in the $q$-adically complete ring $K[[q]]$.
For $j\ge0$, define
\begin{equation}\label{eq:formal-weights}
 u_j=\frac{q^j}{\prod_{\ell=1}^{j}(1-q^\ell)},\qquad
 b_j=-\frac{\alpha q^j}{\prod_{\ell=0}^{j}(1+\alpha q^\ell)}.
\end{equation}
The sums $U=\sum_{j\ge0}u_j$ and $\sum_{j\ge0}b_j$ converge formally.
Moreover, $U(0)=1$, so $U$ is invertible. Put
\begin{equation}\label{eq:a-definition}
 a(q)=\frac{\mu-1-\sum_{j\ge0}b_j}{U},\qquad a_j=a(q)u_j.
\end{equation}
Then
\begin{equation}\label{eq:a0}
 a(0)=\mu-\frac1{1+\alpha},\qquad
 \ordq a_j=\ordq b_j=j.
\end{equation}
Here $1+\alpha$ is nonzero in $K$; the representation is used generically,
not as a direct formula at $\alpha=-1$.

\begin{lemma}\label{lem:formal-moment}
For every $k\ge0$,
\begin{equation}\label{eq:formal-moment}
 v_k=\sum_{j\ge0}\left[a_j(-q^j)^k+b_j(\alpha q^j)^k\right]
 \quad\text{in }K[[q]].
\end{equation}
\end{lemma}
\begin{proof}
Each summand has order at least $j(k+1)$, so the expression is well defined.
Its value at $k=0$ is $\mu-1$ by \eqref{eq:a-definition}.
The coefficient identities
\begin{align*}
 a_j(1-q^j)&=qa_{j-1} &&(j\ge1),\\
 b_j(1+\alpha q^j)&=qb_{j-1} &&(j\ge1),\\
 b_0(1+\alpha)&=-\alpha
\end{align*}
show, on shifting the summation index, that the expression on the right of
\eqref{eq:formal-moment} satisfies
$v_k+v_{k-1}=q^kv_{k-1}-\alpha^k$. For the first family, the $j=0$
contribution to $v_k+v_{k-1}$ vanishes; for the second family it is
$-\alpha^k$. Thus the initial value and recurrence agree with
\eqref{eq:recurrence}, proving the identity.
\end{proof}

There are two nodes at each level $j$:
\begin{equation}\label{eq:formal-nodes}
 x_{j,-}=-q^j,\qquad x_{j,+}=\alpha q^j,
\end{equation}
with weights $a_j$ and $b_j$. If two distinct nodes have levels $j\le k$,
then their difference has order $j$. At the same level this follows from
$x_{j,+}-x_{j,-}=(1+\alpha)q^j$; at different levels it follows from the
nonzero coefficient of the smaller power of $q$.

\subsection{Minimal configurations and the generic order}

\begin{proof}[Proof of \eqref{eq:exact-order}]
Apply the finite Cauchy--Binet identity to truncations of
\eqref{eq:formal-moment}, and then pass coefficientwise to the formal limit.
For a subset of $n$ nodes, its contribution is
\begin{equation}\label{eq:formal-CB}
 \prod_{i=1}^{n}w_ix_i^r\prod_{i<j}(x_j-x_i)^2.
\end{equation}
Write the selected levels in increasing order as
$\ell_1\le\cdots\le\ell_n$. A level occurs at most twice, and the order
of \eqref{eq:formal-CB} is exactly
\begin{equation}\label{eq:level-cost}
 \sum_{i=1}^{n}\bigl[r+2(n-i)+1\bigr]\ell_i.
\end{equation}
All coefficients in this level cost are positive. Every allowable list satisfies
$\ell_i\ge\lfloor(i-1)/2\rfloor$, and equality in each position is
attainable. This is therefore the unique minimizing list of levels.
The formal limit is legitimate: only finitely many node subsets can have
order below a prescribed bound.

The elementary sums
\begin{align*}
 \sum_{i=1}^{n}\fl{\frac{i-1}{2}}
 &=\fl{\frac{(n-1)^2}{4}},\\
 \sum_{i=1}^{n}[2(n-i)+1]\fl{\frac{i-1}{2}}
 &=\binom n3
\end{align*}
give the asserted lower bound. If $n$ is even, both nodes are selected at
every occupied level, so the minimizing subset is unique and its leading
coefficient is nonzero. If $n$ is odd, the two minimizing subsets differ
only in the node chosen at the last level. Their coefficients have different
degrees in $a(0)$, an affine indeterminate in $\mu$, and hence cannot cancel
identically. This proves the exact generic order.
Since $V_n^{(r)}\in R$, the vanishing of its coefficients below this order
also proves divisibility in $R$.
\end{proof}

\subsection{Lowest coefficients and exceptional specializations}

Let
\begin{equation}\label{eq:L-w}
 L=(1+\alpha)\mu-1,\qquad
 w_r=v_r|_{q=0}=\frac{(-1)^rL-\alpha^{r+1}}{1+\alpha}.
\end{equation}
The expression for $w_r$ is a polynomial after cancellation of the apparent
denominator; alternatively it is determined by $w_0=\mu-1$ and
$w_r=-w_{r-1}-\alpha^r$.

\begin{proposition}\label{prop:leading}
Writing $e=e_{n,r}$, the lowest coefficients are
\begin{align}
 [q^e]V_{2m}^{(r)}
 &=(-1)^{m(r+1)}\alpha^{m(2m+r-1)}L^m,\label{eq:lead-even}\\
 [q^e]V_{2m+1}^{(r)}
 &=(-1)^{m(r+1)}\alpha^{m(2m+r+1)}L^m w_r.\label{eq:lead-odd}
\end{align}
Formula \eqref{eq:lead-even} applies for $m\ge1$ and
\eqref{eq:lead-odd} for $m\ge0$.
\end{proposition}
\begin{proof}
At every occupied level, the leading weight coefficients are
$A=L/(1+\alpha)$ and $B=-\alpha/(1+\alpha)$.
For $n=2m$, the product of all weight coefficients is $(AB)^m$.
The factors $x_i^r$ contribute $(-1)^{mr}\alpha^{mr}$.
Each same-level Vandermonde square contributes $(1+\alpha)^2$,
cancelling the corresponding denominator in $AB$.
For a positive node at level $j$, the higher levels contribute a further
factor $\alpha^{4(m-1-j)}$ to the leading Vandermonde coefficient.
Multiplication gives \eqref{eq:lead-even}.

For $n=2m+1$, both nodes at each level $0,\ldots,m-1$ are selected.
The extra node contributes an additional $\alpha^{2m}$ through its
interactions with the preceding positive nodes. The sum of its two possible
weight-and-shift coefficients is
$A(-1)^r+B\alpha^r=w_r$.
This yields \eqref{eq:lead-odd}.
\end{proof}

\begin{remark}\label{rem:exception}
For $\alpha=2,\mu=3,r=2$, one has $w_2=0$. Thus at $n=5$ the
coefficient of the generic lowest power $q^{18}$ vanishes.
Exact polynomial calculation gives order $19$ in this example
(\cref{app:checks}). This is an exceptional specialization, not a
failure of the generic identity. No equality after every specialization is
claimed here.
\end{remark}

\subsection{The highest degree}

\begin{proof}[Proof of \eqref{eq:exact-degree}]
Induction in \eqref{eq:recurrence} gives
$\degq v_k=k(k+1)/2$ with leading coefficient $\mu-1$.
For a permutation $\pi$ of $\{0,\ldots,n-1\}$, the corresponding
term of the determinant has degree
$\sum_i(i+\pi(i)+r)(i+\pi(i)+r+1)/2$.
The only part depending on the permutation is $\sum_i i\pi(i)$.
Indeed, the difference between its value at the identity and at $\pi$ is
$\frac12\sum_i(i-\pi(i))^2$, which is positive unless $\pi$ is the identity.
Thus the identity permutation uniquely supplies the highest power of $q$,
whose coefficient is $(\mu-1)^n$. This proves the degree formula.
\end{proof}
For fixed $t\ge0$ and $r=\lfloor tn\rfloor$, these algebraic identities give
\begin{equation}\label{eq:algebra-rates}
 \frac{D_{n,r}}{n^3}\longrightarrow C(t),\qquad
 \frac{e_{n,r}}{n^3}\longrightarrow b_0(t):=\frac16+\frac t4.
\end{equation}

\section{The prescribed resonant cyclotomic divisor}\label{sec:cyclotomic}

This section gives a self-contained construction of the prescribed divisor.
The proof establishes a guaranteed divisor; it does not establish that all
its cyclotomic exponents are exact.

Fix an integer $d\ge1$ and put $h=q^d-1$, $a=\alpha^d$. Let $\mathsf T_d$
denote the shift by $d$ on sequences and define
\begin{equation}\label{eq:W}
 W_{s,k}=(\mathsf T_d-a)^s v_k
 =\sum_{j=0}^{s}\binom sj(-a)^{s-j}v_{k+jd}.
\end{equation}
We use the sans-serif symbol $\mathsf T_d$ for this operator, to distinguish
it from the scalar sum $T_d(n,p)$ in \eqref{eq:T}.

\begin{lemma}[Resonance and lifting]\label{lem:lifting}
For $s\ge1$ and $k\ge1$,
\begin{equation}\label{eq:W-rec}
 W_{s,k}=(q^k-1)W_{s,k-1}
 +q^k\sum_{j=1}^{s}\binom sj h^jW_{s-j,k+jd-1}.
\end{equation}
For all $s,k\ge0$,
\begin{equation}\label{eq:W-div}
 h^{\min(s,\lfloor k/d\rfloor)}\mid W_{s,k}\quad\text{in }R.
\end{equation}
\end{lemma}
\begin{proof}
Substitution of \eqref{eq:recurrence} into \eqref{eq:W} cancels the
inhomogeneous terms because their sum is
$-\alpha^k\sum_{j=0}^s\binom sj(-a)^{s-j}a^j=0$.
Writing $q^{k+jd}-1=(q^k-1)+q^k((1+h)^j-1)$, the remaining terms
are reorganized by the commuting-operator identity
\[
 ((1+h)\mathsf T_d-a)^s=((\mathsf T_d-a)+h\mathsf T_d)^s.
\]
This gives \eqref{eq:W-rec}.

For \eqref{eq:W-div}, use induction on $s$, and then on $k$ for fixed $s$.
The claim is immediate when $s=0$ or $k<d$. In the first term of
\eqref{eq:W-rec}, the floor $\lfloor k/d\rfloor$ either stays unchanged
from $k-1$, or increases by one when $d\mid k$; in the latter case
$h\mid q^k-1$ supplies the additional factor.
In a term with $j\ge1$, the explicit power of $h$ and the induction
hypothesis give an exponent at least
\[
 j+\min\left(s-j,\fl{\frac{k+jd-1}{d}}\right)
 \ge\min\left(s,\fl{\frac{k}{d}}+2j-1\right)
 \ge\min\left(s,\fl{\frac{k}{d}}\right).
\]
This completes both inductions.
\end{proof}

\begin{proposition}[Shifted divisor]\label{prop:divisor}
The polynomial $\DD_{n,r}$ of \eqref{eq:divisor} divides $V_n^{(r)}$ in $R$.
For $r=0$ it can also be written as
\begin{equation}\label{eq:unshifted-divisor}
 \DD_{n,0}(q)=q^{\binom n3}
 \prod_{\ell\ge1}(q^\ell-1)^{2(n-2\ell)_+}.
\end{equation}
\end{proposition}
\begin{proof}
Let integers $a_i,b_j$ satisfy
\begin{equation}\label{eq:admissible}
 0\le a_i\le\fl{\frac i d},\qquad
 0\le b_j\le\fl{\frac j d},\qquad
 a_i+b_j\le\fl{\frac{i+j+r}{2d}}.
\end{equation}
Apply lower unitriangular row transformations
\[
 R_i\longmapsto\sum_{u=0}^{a_i}\binom{a_i}{u}
 (-\alpha^d)^{a_i-u}R_{i-(a_i-u)d}
\]
and analogous column transformations with orders $b_j$.
All indices are nonnegative by \eqref{eq:admissible}. The transformed
entry is
\[
 W_{a_i+b_j,k_0},\qquad k_0=i+j+r-d(a_i+b_j)\ge d(a_i+b_j).
\]
By \cref{lem:lifting}, this entry is divisible by $h^{a_i+b_j}$.
Factoring $h^{a_i}$ from row $i$ and $h^{b_j}$ from column $j$ gives
$h^{\sum_i a_i+\sum_j b_j}\mid V_n^{(r)}$.

For $p+p'=r$, choose
\[
 a_i=\min\left\{\fl{\frac i d},\fl{\frac{i+p}{2d}}\right\},\qquad
 b_j=\min\left\{\fl{\frac j d},\fl{\frac{j+p'}{2d}}\right\}.
\]
The floor inequality $\lfloor x\rfloor+\lfloor y\rfloor\le\lfloor x+y\rfloor$
verifies \eqref{eq:admissible}. Thus
$(q^d-1)^{T_d(n,p)+T_d(n,p')}\mid V_n^{(r)}$.
In particular, $\Phi_d(q)^{\nu_d(n,r)}\mid V_n^{(r)}$.
The cyclotomic polynomials are distinct monic irreducibles over $\Q$,
so their product divides over $\Q[\alpha,\mu][q]$ and, by monic division,
over $R$. It is coprime to $q$, and \cref{thm:main-algebra} supplies the
remaining $q$-power.

At $r=0$, $\nu_d(n,0)=2S_d(n)$ and
$S_d(n)=\sum_{m\ge1}(n-2dm)_+$.
The exponent of $\Phi_d(q)$ in the product on the right of
\eqref{eq:unshifted-divisor} is therefore exactly $2S_d(n)$.
\end{proof}

In particular,
\begin{equation}\label{eq:unshifted-degree}
 \degq\DD_{n,0}=
 \begin{cases}
 n(n-1)^2/4,&n\text{ odd},\\
 n^2(n-2)/4,&n\text{ even},
 \end{cases}
 \qquad \degq\DD_{n,0}=\frac14n^3+O(n^2).
\end{equation}
This is the divisor form of the expectation recorded in
\cite[Remark 7]{KRVZ}. The factors $(q^\ell-1)$ in
\eqref{eq:unshifted-divisor} are not themselves pairwise coprime;
reassembly through their irreducible cyclotomic factors is essential.

\section{Positive moments and sharp analytic bounds}\label{sec:positive}

\subsection{The product and the actual remainder}

Comparison of coefficients in \eqref{eq:exponential} gives
\begin{equation}\label{eq:functional}
 E_q(qz)=(1+z)E_q(z).
\end{equation}
Iterating downward and using $E_q(0)=1$ yields, for $|q|>1$,
\begin{equation}\label{eq:product}
 E_q(z)=\prod_{j\ge1}(1+zq^{-j}).
\end{equation}
The product converges locally uniformly in $z$. In the remainder of this
section assume $q>1,\alpha>0$ and $\mu=E_q(\alpha)$.

\begin{proof}[Proof of the moment and positivity assertions in
\cref{thm:main-positive}]
Let $P_j=\prod_{\ell=1}^j(1+\alpha q^{-\ell})$ and $P_0=1$.
The weights in \eqref{eq:positive-nodes} satisfy $w_j=P_j-P_{j-1}$,
so telescoping gives $\sum_{j\ge1}w_j=\mu-1=v_0$.
They also satisfy
\[
 w_1q=\alpha,\qquad w_{j+1}q^{j+1}-w_jq^j=\alpha w_j.
\]
Since $0<w_j\le\alpha\mu q^{-j}$, all moment sums converge absolutely.
The displayed relations, followed by an index shift, verify the recurrence
for $\sum_{j\ge1}w_j(\alpha q^{-j})^k$.
This proves the moment identity by uniqueness of the recurrence solution.

Finite Cauchy--Binet followed by the convergent limit now gives
\begin{equation}\label{eq:positive-CB}
 V_n^{(r)}=
 \sum_{1\le j_1<\cdots<j_n}
 \prod_{a=1}^n w_{j_a}x_{j_a}^r
 \prod_{a<b}(x_{j_b}-x_{j_a})^2.
\end{equation}
Each term is positive, and at least one exists. The finite determinants
converge because their entries do. This proves strict positivity and the
convergence of the nonnegative expansion.
\end{proof}

\subsection{A two-sided product estimate}

\begin{proof}[Proof of \eqref{eq:main-bounds} and \eqref{eq:log-asymptotic}]
Retain the subset $j_a=a$ in \eqref{eq:positive-CB} and use
$w_a\ge\alpha q^{-a}$. Since
\[
 x_{j_a}-x_{j_b}=\alpha q^{-j_a}(1-q^{-(j_b-j_a)})\quad(a<b),
\]
this yields the lower bound in \eqref{eq:main-bounds}; its $q$-exponent is
\[
 \sum_{a=1}^n[r+2(n-a)+1]a=K_{n,r}.
\]
For the upper bound, use $w_j\le\alpha\mu q^{-j}$ and remove all the
factors $(1-q^{-(j_b-j_a)})^2\le1$. It remains to sum
\begin{equation}\label{eq:gap-sum}
 \sum_{j_1<\cdots<j_n}
 q^{-\sum_{a=1}^n[r+2(n-a)+1]j_a}
 =q^{-K_{n,r}}\prod_{m=1}^n(1-q^{-m(m+r)})^{-1}.
\end{equation}
Indeed, write $j_a=a+g_1+\cdots+g_a$, where each $g_a\ge0$.
The coefficient of $g_a$ in the exponent is
$(n-a+1)(n-a+1+r)$, and the $n$ sums are geometric.
The total power of $\alpha$ is $n+nr+n(n-1)=n^2+rn$.
This proves both bounds.

The lower product is at least
$\prod_{k\ge1}(1-q^{-k})^{2n}$.
The upper reciprocal product is at most
$\prod_{m\ge1}(1-q^{-m^2})^{-1}$, independent of $r$ and $n$.
Taking logarithms proves the uniform error in \eqref{eq:log-asymptotic}.
For $r=\lfloor tn\rfloor$, division by $n^3\log q$ gives
$A(t)=1/3+t/2$.
\end{proof}

\begin{corollary}[No faster cubic decay for the same positive determinant]
For fixed $q>1,\alpha>0$ at $\mu=E_q(\alpha)$ and fixed $t\ge0$,
\[
 \lim_{n\to\infty}\frac{-\log V_n^{(\lfloor tn\rfloor)}}{n^3\log q}
 =\frac13+\frac t2.
\]
Consequently no valid eventual estimate for this same sequence can replace
this exponent by a larger fixed constant in a bound of the form
$q^{-An^3+o(n^3)}$.
\end{corollary}

This conclusion uses positivity and two-sided bounds. It is not an inference
from a one-sided determinant estimate, and it is not asserted for arbitrary
complex specializations.

\section{The divisor-degree budget}\label{sec:budget}

\begin{lemma}[A floor-sum inequality]\label{lem:floor}
For integers $L,n\ge1$, define $G(p)=\sum_{i=0}^{n-1}\lfloor(i+p)/L\rfloor$
for $p\ge0$. Then
\begin{equation}\label{eq:G-super}
 G(p)+G(p')\le G(0)+G(p+p')\quad(p,p'\ge0).
\end{equation}
Consequently,
\begin{equation}\label{eq:nu-budget}
 \nu_d(n,r)\le S_d(n)+S_d(n+r)-S_d(r).
\end{equation}
\end{lemma}
\begin{proof}
Write $n=aL+b$ with $0\le b<L$, and $p=cL+s$ with $0\le s<L$.
Direct summation gives
\[
 G(p)-G(0)=ap+bc+(b+s-L)_+.
\]
The last two terms form
$H(p)=b\lfloor p/L\rfloor+(b+(p\bmod L)-L)_+$.
Consider the periodic binary sequence with $L-b$ consecutive zeros followed
by $b$ consecutive ones in each period. The minimum number of ones in an
interval of length $p$ is exactly $H(p)$: complete periods contribute $bc$,
and a residual interval of length $s$ contributes at least
$(s-(L-b))_+$, a minimum attainable by placing it across the zero block.
Splitting an arbitrary interval of length $p+p'$ into pieces of lengths
$p,p'$ shows that its number of ones is at least $H(p)+H(p')$.
Taking a minimum proves superadditivity of $H$, hence \eqref{eq:G-super}.

Take $L=2d$. Since $T_d(n,p)\le G(p)$, for $p+p'=r$ one obtains
\[
 T_d(n,p)+T_d(n,p')\le G(0)+G(r)
 =S_d(n)+S_d(n+r)-S_d(r).
\]
Maximization over the split proves \eqref{eq:nu-budget}.
\end{proof}

\begin{proposition}\label{prop:budget}
For every $n,r$,
\begin{equation}\label{eq:exact-budget}
 \degq\DD_{n,r}\le e_{n,r}+F(n)+F(n+r)-F(r),
\end{equation}
where
\begin{equation}\label{eq:F}
 F(N):=\sum_{d\ge1}\varphi(d)S_d(N)
 =\sum_{\ell\ge1}\ell(N-2\ell)_+
 =\frac{N^3}{24}+O(N^2).
\end{equation}
In particular, for $r=\lfloor tn\rfloor$ and fixed $t\ge0$,
\begin{equation}\label{eq:budget-limsup}
 \limsup_{n\to\infty}\frac{\degq\DD_{n,r}}{n^3}
 \le\frac14+\frac{3t}{8}+\frac{t^2}{8}.
\end{equation}
\end{proposition}
\begin{proof}
Use $S_d(N)=\sum_{m\ge1}(N-2dm)_+$ and
$\sum_{d\mid\ell}\varphi(d)=\ell$ to obtain the middle expression in
\eqref{eq:F}. Summing the resulting quadratic polynomial in $\ell$ gives its
asymptotic. Multiply \eqref{eq:nu-budget} by $\varphi(d)$ and sum over $d$
to prove \eqref{eq:exact-budget}. Since
\[
 \frac{1+(1+t)^3-t^3}{24}
 =\frac1{12}+\frac t8+\frac{t^2}{8},
\]
combination with $b_0(t)=1/6+t/4$ proves \eqref{eq:budget-limsup}.
\end{proof}

\begin{remark}[Which divisor is bounded]
The right side of \eqref{eq:exact-budget} bounds the degree of the
\emph{specified} polynomial \eqref{eq:divisor}. It does not bound the degree
of every polynomial dividing $V_n^{(r)}$, nor the full cyclotomic valuation
of a specialized determinant. Replacing this distinction by an assertion of
universal sharpness would not be justified by the proof.
\end{remark}

\section{Arithmetic comparison and shift optimality}\label{sec:arithmetic}

For positive rational $q=\rho/\sigma>1$ in lowest terms, with
$\rho>\sigma\ge1$, define
\begin{equation}\label{eq:gamma}
 \gamma_q=\frac{\log\rho}{\log\sigma}\quad(\sigma>1),\qquad
 \gamma_q=\infty\quad(\sigma=1).
\end{equation}
The symbol $\gamma_q$ is a height ratio; it is not Euler's constant.
We give the norm calculation explicitly to specify the comparison whose
optimality is at issue. It is the shifted counterpart of
\cite[Section 5]{KRVZ}.

\subsection{Denominator clearing and conjugates}

For fixed $q>1$, the cyclotomic product satisfies
\begin{equation}\label{eq:divisor-size}
 \log|\DD_{n,r}(q)|=(\degq\DD_{n,r})\log q+O_q(n^2\log(n+1)).
\end{equation}
To see this, use the M\"obius product formula to bound
$|\log|\Phi_d(q)|-\varphi(d)\log q|$ by
$\sum_{m\ge1}|\log(1-q^{-m})|$, a finite constant independent of $d$.
Also $\nu_d(n,r)\le2\sum_{i<n}\lfloor i/d\rfloor\le n^2/d$, and
$\nu_d(n,r)=0$ for $d\ge n$. Summing proves \eqref{eq:divisor-size}.

Suppose that $\alpha>0$ and $\mu=E_q(\alpha)$ are algebraic, and put
$K=\Q(\alpha,\mu)$, $D=[K:\Q]$. Choose an integer $c\ge1$ such that
$c\alpha$ and $c\mu$ are algebraic integers. The polynomial bounds
\[
 \deg_\alpha V_n^{(r)}\le n(n-1+r),\qquad
 \deg_\mu V_n^{(r)}\le n
\]
follow from the recurrence and determinant expansion. They remain valid
for the quotient by the monic $q$-only polynomial $\DD_{n,r}$.
Writing $b_{n,r}=\degq\DD_{n,r}$, it follows that
\begin{equation}\label{eq:integer}
 I_{n,r}:=c^{n(n+r)}\sigma^{D_{n,r}-b_{n,r}}
 \frac{V_n^{(r)}(q,\alpha,\mu)}{\DD_{n,r}(q)}
 \quad\text{is an algebraic integer.}
\end{equation}
It is nonzero by \cref{thm:main-positive}.

For any embedding $\iota:K\hookrightarrow\C$, solving the recurrence as in
\eqref{eq:tail}, without imposing the actual-value identity on
$\iota(\mu)$, gives a bound
\[
 |\iota(v_k)|\le C_\iota q^{k(k+1)/2}.
\]
Indeed, the corresponding finite partial sums of
$\sum\iota(\alpha)^j/\prod_{\ell=1}^j(q^\ell-1)$ are bounded.
The determinant expansion and the same degree maximization as in
\cref{thm:main-algebra} therefore give
\begin{equation}\label{eq:conjugates}
 |\iota(V_n^{(r)})|\le n!\,C_\iota^n q^{D_{n,r}}.
\end{equation}
At the distinguished real embedding the much stronger estimate
\eqref{eq:log-asymptotic} applies.

\begin{proposition}[The cubic norm criterion]\label{prop:norm}
Fix $t\ge0$ and $r=\lfloor tn\rfloor$. Along any infinite subsequence on
which $b_{n,r}/n^3\to B$, the assumption $[\Q(\alpha,\mu):\Q]=D$ is
contradicted if
\begin{equation}\label{eq:norm-condition}
 A(t)+C(t)-D(C(t)-B)>0,
 \qquad
 \gamma_q>\frac{A(t)+C(t)}{A(t)+C(t)-D(C(t)-B)}.
\end{equation}
For integer $q$, the second condition is automatic when the first holds.
\end{proposition}
\begin{proof}
The logarithm of the $c$-factor in \eqref{eq:integer} is $O(n^2)$ because
$r=O(n)$. Multiply the distinguished bound for $I_{n,r}$ by the
$D-1$ bounds arising from \eqref{eq:conjugates}, and use
\eqref{eq:divisor-size}. The result is
\begin{align}\label{eq:norm-exponent}
 \log|N_{K/\Q}(I_{n,r})|
 &\le\bigl((A+C)\log\sigma
       -(A+C-D(C-B))\log\rho\bigr)n^3+o(n^3).
\end{align}
The coefficient is strictly negative under \eqref{eq:norm-condition}.
But the norm is a nonzero integer, whose absolute value is at least one.
This is a contradiction along the subsequence.
\end{proof}

\begin{corollary}[Unshifted consequences in the positive range]\label{cor:arithmetic}
Let $q=\rho/\sigma>1$ be rational and $\alpha>0$ algebraic.
If $\gamma_q>6$, then $\alpha$ and $E_q(\alpha)$ cannot both lie in a
number field of degree at most two. In particular, for positive rational
$\alpha$, $E_q(\alpha)$ is neither rational nor quadratic algebraic.
If $\alpha>0$ is rational and $\gamma_q>12/7$, then $E_q(\alpha)$ is
irrational.
\end{corollary}
\begin{proof}
At $t=0$, \eqref{eq:unshifted-degree} gives $B=1/4$ and
$A=1/3,C=2/3$. Substitution of $D=1,2$ into
\eqref{eq:norm-condition} gives $12/7$ and $6$, respectively.
\end{proof}
These known arithmetic values agree with \cite[Corollary 1.1]{Rochev2011}.
The positive-range corollary is proved directly here to identify exactly
which norm comparison is subject to the shift barrier. The more general
complex-parameter results in \cite{KRVZ,Rochev2011} are not needed for
this positivity argument.

\subsection{The global shift barrier}

\begin{proof}[Proof of \cref{thm:main-optimality}]
Divisibility is \cref{prop:divisor}; the constants $A,C$ follow from
\cref{thm:main-algebra,thm:main-positive}; and \cref{prop:budget} gives
$B\le\overline B(t)$. Since $A+B>0$ and the ratio decreases with $B$,
\[
 \frac{A+C}{A+B}\ge\frac{A+C}{A+\overline B(t)}
 =\frac{24+36t+12t^2}{14+21t+3t^2}.
\]
The exact difference from the unshifted value is
\begin{equation}\label{eq:strict-gap}
 \frac{24+36t+12t^2}{14+21t+3t^2}-\frac{12}{7}
 =\frac{48t^2}{7(14+21t+3t^2)}.
\end{equation}
This is nonnegative, and positive for $t>0$.
For degree two, set $\eta=(A+B)/(A+C)$. The preceding inequality implies
$\eta\le7/12$. When $2\eta-1>0$,
\[
 \tau_2=\frac1{2\eta-1}\ge6.
\]
Otherwise the cubic comparison gives no finite degree-two critical value.
Finally, \eqref{eq:unshifted-degree} shows that the bounds are attained at
$t=0$. This proves the theorem, including its accumulation-point formulation.
\end{proof}

\begin{remark}[What the theorem does not rule out]\label{rem:scope}
The exact analytic exponent is established for $q>1,\alpha>0$ at the actual
value, and the barrier uses precisely the divisor \eqref{eq:divisor} and
the conjugate bound \eqref{eq:conjugates}. It is not a universal theorem
about irrationality thresholds. In particular, it leaves open extra
specialization-dependent content, a larger uniform divisor, different
minors or determinants, other conjugate bounds, and other denominator
clearings. It neither determines every shifted cyclotomic valuation nor
excludes a future general-parameter result below $12/7$.
\end{remark}

\section{Entire normalization and exact forced zeros}
\label{sec:entire-normalization}

In this section and the subsequent analytic sections, fix a real number
$q>1$ and put $p=q^{-1}$. The parameter $\alpha$ is now allowed to range
over $\C$. Write
\begin{equation}\label{eq:analytic-products}
 \pi_j=(p;p)_j=\prod_{\ell=1}^j(1-p^\ell),\qquad
 P_\infty=(p;p)_\infty,\qquad
 \mathcal M=\prod_{\ell\ge1}(1-p^\ell)^{-\ell}.
\end{equation}
As usual, $(a;p)_j=\prod_{\ell=0}^{j-1}(1-ap^\ell)$, and an empty
product is one. Set
\begin{align}
 H_{n,r}(\alpha)&=V_n^{(r)}(q,\alpha,E_q(\alpha)),\label{eq:analytic-H}\\
 \mathcal A_{n,r}
 &=\prod_{i,j=0}^{n-1}(1-p^{r+i+j+1})^{-1},\label{eq:analytic-A}\\
 \mathcal C_{n,r}
 &=\det\left(\frac{p^{r+i+j+1}}{1-p^{r+i+j+1}}\right)_{i,j=0}^{n-1}
 =p^{K_{n,r}}\prod_{\ell=1}^{n-1}(1-p^\ell)^{2(n-\ell)}
   \mathcal A_{n,r},\label{eq:analytic-C}\\
 \mathcal L_{n,r}(\alpha)
 &=\alpha^{n^2+rn}
   \prod_{m=1}^{n-1}(1+\alpha p^m)^{n-m},\label{eq:analytic-L}
\end{align}
where $K_{n,r}$ is as in \cref{thm:main-positive}.
The evaluation in \eqref{eq:analytic-C} is the Cauchy determinant formula;
in particular, $\mathcal C_{n,r}>0$.

\begin{theorem}[Exact forced zeros]\label{thm:forced-alpha-orders}
For $n\ge1$ and $r\ge0$, the function $H_{n,r}$ is entire and
\begin{equation}\label{eq:forced-alpha-orders}
 \operatorname{ord}_{\alpha=0}H_{n,r}=n^2+rn,\qquad
 \operatorname{ord}_{\alpha=-p^{-m}}H_{n,r}=n-m
 \quad(1\le m<n).
\end{equation}
Consequently, the quotient
\begin{equation}\label{eq:analytic-normalization}
 \Psi_{n,r}(\alpha)
 =\frac{H_{n,r}(\alpha)}
 {\mathcal L_{n,r}(\alpha)\mathcal C_{n,r}}
\end{equation}
extends to an entire function, and $\Psi_{n,r}(0)=1$.
\end{theorem}

\begin{proof}
The recurrence defining $v_k$ shows directly that $H_{n,r}$ is entire.
The moment representation from \cref{sec:positive} extends to complex
$\alpha$:
\begin{equation}\label{eq:complex-moment-cb}
 \begin{split}
 v_k&=\sum_{j\ge1}w_jx_j^k,\qquad
 x_j=\alpha p^j,\qquad
 w_j=\alpha p^j\prod_{\ell=1}^{j-1}(1+\alpha p^\ell),\\
 H_{n,r}(\alpha)
 &=\sum_{1\le j_1<\cdots<j_n}
 \prod_{a=1}^n w_{j_a}x_{j_a}^{r}
 \prod_{a<b}(x_{j_b}-x_{j_a})^2.
 \end{split}
\end{equation}
Indeed, both series converge normally on compact subsets of $\C$.
For example, on $|\alpha|\le R$ the partial products in $w_j$ are at
most $\exp(Rp/(1-p))$ in absolute value. The remaining geometric factors
give an absolutely summable bound for the Cauchy--Binet series.
The moment identity follows either from the recurrence or from the
identity theorem applied to the positive real axis.

At $\alpha=0$, extracting the powers of $\alpha$ in
\eqref{eq:complex-moment-cb} gives
\[
 H_{n,r}(\alpha)
 =\alpha^{n^2+rn}\bigl(\mathcal C_{n,r}+O(\alpha)\bigr).
\]
Since $\mathcal C_{n,r}>0$, this proves the first assertion of
\eqref{eq:forced-alpha-orders} and the value of $\Psi_{n,r}(0)$.

Fix $m<n$ and let $\alpha_0=-p^{-m}$. The first $m$ weights in
\eqref{eq:complex-moment-cb} are nonzero at $\alpha_0$; every subsequent
weight has a simple zero there. Thus every summand has order at least
$n-m$. Terms of exactly this order must contain all of the first $m$
nodes. There is no cancellation among their leading coefficients:
\[
 \operatorname{sgn}w_j(\alpha_0)=(-1)^j\quad(j\le m),\qquad
 \operatorname{sgn}w_j'(\alpha_0)=(-1)^m\quad(j>m).
\]
The factors $x_j^r$ contribute the common sign $(-1)^{rn}$, and the
Vandermonde squares are strictly positive. At least one such summand
exists. Normal convergence permits coefficient extraction, so the sum
of these nonzero coefficients is nonzero. This proves the second
assertion and removes every apparent pole in
\eqref{eq:analytic-normalization}.
\end{proof}

\section{A Schur expansion and a uniform scaling limit}
\label{sec:schur-edge-limit}

We use the standard Schur identities in \cite[Chapter~I]{Macdonald}.
For a partition $\lambda$, let $|\lambda|$ and $\ell(\lambda)$ denote
its size and length. A cell $\square=(i,b)$ in its diagram has content
$c(\square)=b-i$ and hook length $h(\square)$, and
$b(\lambda)=\sum_i(i-1)\lambda_i$. Put
\[
 X_n=(1,p,\ldots,p^{n-1}),\qquad X=(1,p,p^2,\ldots).
\]
The principal specializations are
\begin{equation}\label{eq:principal-schur-specialization}
 s_\lambda(X)=
 \frac{p^{b(\lambda)}}{\prod_{\square\in\lambda}(1-p^{h(\square)})},
 \qquad
 s_\lambda(X_n)=s_\lambda(X)
 \prod_{\square\in\lambda}(1-p^{n+c(\square)})
 \quad(\ell(\lambda)\le n),
\end{equation}
and $s_\lambda(X_n)=0$ for $\ell(\lambda)>n$.

\begin{proposition}[Exact partition expansion]\label{prop:schur-representation}
Define
\begin{equation}\label{eq:scaled-F-definition}
 F_{n,r}(z)=\Psi_{n,r}(p^{-n-r}z),\qquad a=p^r.
\end{equation}
Then, for every $z\in\C$,
\begin{equation}\label{eq:scaled-schur-sum}
 F_{n,r}(z)=\frac{1}{\mathcal A_{n,r}}
 \sum_{\lambda}p^{|\lambda|}s_\lambda(X_n)^2
 \prod_{\square\in\lambda}\bigl(a+zp^{c(\square)}\bigr).
\end{equation}
The series converges normally on compact subsets of $\C$. Its
normalizing factor satisfies
\begin{equation}\label{eq:schur-partition-normalizer}
 \mathcal A_{n,r}
 =\sum_\lambda p^{(r+1)|\lambda|}s_\lambda(X_n)^2.
\end{equation}
\end{proposition}

\begin{proof}
For a configuration $j_1<\cdots<j_n$ in
\eqref{eq:complex-moment-cb}, write
$j_a=a+\lambda_{n-a+1}$. The alternant definition of Schur polynomials
gives the Vandermonde ratio $s_\lambda(X_n)$. Relative to the ground
configuration $j_a=a$, the geometric weights and the squared
Vandermonde therefore contribute
$p^{(r+1)|\lambda|}s_\lambda(X_n)^2$. The remaining product is
\begin{equation}\label{eq:configuration-extra-product}
 \prod_{a=1}^n\prod_{\ell=a}^{j_a-1}(1+\alpha p^\ell)
 =\prod_{\square\in\lambda}(1+\alpha p^{n+c(\square)}).
\end{equation}
Ordinary Cauchy summation gives
\eqref{eq:schur-partition-normalizer}. Substitution of
$\alpha=p^{-n-r}z$ in \eqref{eq:configuration-extra-product}, together
with $p^{r|\lambda|}$ from the weights, gives
\eqref{eq:scaled-schur-sum}. Normal convergence also follows from the
summable majorant established in the next lemma.
\end{proof}

For $0\le a\le1$, define
\begin{equation}\label{eq:infinite-normalizer-and-Phi}
 \mathcal A_\infty(a)=\prod_{\ell\ge1}(1-ap^\ell)^{-\ell},
 \qquad
 \Phi_p(z)=\prod_{\ell\ge1}(1+zp^\ell)^\ell.
\end{equation}
The second product is entire, and $\mathcal A_\infty(1)=\mathcal M$.
We also write
\begin{equation}\label{eq:infinite-shifted-normalizer}
 \mathcal A_{\infty,r}=\mathcal A_\infty(p^r)
 =\prod_{\ell\ge1}(1-p^{r+\ell})^{-\ell}.
\end{equation}

\begin{lemma}[Summation and a content moment]\label{lem:super-schur-bound}
For $0\le a\le1$ and $z\in\C$,
\begin{equation}\label{eq:infinite-super-cauchy}
 \sum_\lambda p^{|\lambda|}s_\lambda(X)^2
 \prod_{\square\in\lambda}(a+zp^{c(\square)})
 =\mathcal A_\infty(a)\Phi_p(z).
\end{equation}
The convergence is locally uniform in $z$ and uniform in $a\in[0,1]$.
Furthermore, if
$C_1(\lambda)=\sum_{\square\in\lambda}p^{c(\square)}$, then for $R\ge0$
\begin{equation}\label{eq:content-moment-bound}
 \sum_\lambda p^{|\lambda|}s_\lambda(X)^2 C_1(\lambda)
 \prod_{\square\in\lambda}(1+Rp^{c(\square)})
 =\mathcal M\Phi_p(R)\frac{(1+R)p}{(1-p)^2}.
\end{equation}
\end{lemma}

\begin{proof}
Use the supersymmetric specialization
$\mathsf A_{a,z}=aX+\omega(zX)$, where $\omega$ is the standard
involution on symmetric functions. Equivalently, its power sums are
\[
 p_k[\mathsf A_{a,z}]
 =\frac{a^k+(-1)^{k-1}z^k}{1-p^k}.
\]
The hook-content specialization gives
\begin{equation}\label{eq:super-hook-content}
 s_\lambda[\mathsf A_{a,z}]
 =s_\lambda(X)\prod_{\square\in\lambda}(a+zp^{c(\square)}).
\end{equation}
For $a>0$ this follows by homogeneity from the usual specialization
$p_k\mapsto(1-u^k)/(1-p^k)$ with $u=-z/a$; the case $a=0$ follows
polynomially.

Here it is important to justify Cauchy summation for arbitrarily large
$|z|$. First take finite alphabets $aX_N$, $RX_N$, and $pX_N$, with
$R\ge0$. The mixed ordinary and dual Cauchy identity gives
\[
 \sum_\lambda s_\lambda[aX_N+\omega(RX_N)]s_\lambda(pX_N)
 =\prod_{i,j=0}^{N-1}\frac{1+Rp^{i+j+1}}{1-ap^{i+j+1}}.
\]
All terms on the left are nonnegative, by the skew-Schur expansion
of a supersymmetric Schur function. Increasing $N$ and using monotone
convergence proves that the corresponding infinite sum equals
$\mathcal A_\infty(a)\Phi_p(R)$, which is finite for every $R$.
For complex $z$ with $|z|\le R$ and $a\in[0,1]$, the absolute value of
each summand in \eqref{eq:infinite-super-cauchy} is bounded by the
summand with $(a,z)=(1,R)$. This proves normal convergence and the
identity by analytic continuation. In particular, no logarithmic power
series with a restricted $z$-radius is being used.

To obtain \eqref{eq:content-moment-bound}, apply the same Cauchy identity
to $\mathsf A_{1,R}$ and $p\mathsf A_{1,\varepsilon}$. For sufficiently
small nonnegative $\varepsilon$ it gives
\[
 \sum_\lambda p^{|\lambda|}s_\lambda(X)^2
 \prod_{\square\in\lambda}(1+Rp^{c(\square)})
                            (1+\varepsilon p^{c(\square)})
 =\mathcal M\Phi_p(R)\Phi_p(\varepsilon)
   \prod_{\ell\ge1}(1-R\varepsilon p^\ell)^{-\ell}.
\]
Taking the coefficient of $\varepsilon$ proves the assertion, since
$\sum_{\ell\ge1}\ell p^\ell=p/(1-p)^2$. Coefficient extraction is
justified either by the preceding finite-alphabet argument and
nonnegativity, or by the convergent identity in a neighborhood of zero.
\end{proof}

\begin{theorem}[Uniform shifted scaling limit]\label{thm:scaled-limit}
For every $R\ge0$, every $n\ge1$, and every integer $r\ge0$,
\begin{equation}\label{eq:scaled-limit-explicit}
 \sup_{|z|\le R}|F_{n,r}(z)-\Phi_p(z)|
 \le
 \frac{2\mathcal M\Phi_p(R)(R+2)}{(1-p)^2}\,p^{n+1}.
\end{equation}
Thus $F_{n,r}\to\Phi_p$ locally uniformly on $\C$, uniformly in the
entire range $r\ge0$. The same conclusion, with an $O_{p,R,d}(p^n)$
bound, holds for every fixed derivative order $d$ on $|z|\le R$.
\end{theorem}

\begin{proof}
Let $t_{\lambda,n}=s_\lambda(X_n)/s_\lambda(X)$. By
\eqref{eq:principal-schur-specialization}, $0\le t_{\lambda,n}\le1$ and
\begin{equation}\label{eq:finite-alphabet-error}
 1-t_{\lambda,n}^2\le 2p^n C_1(\lambda).
\end{equation}
When $\ell(\lambda)\le n$, this is the elementary inequality
$1-\prod(1-u_j)^2\le2\sum u_j$ applied to
$u_j=p^{n+c(\square)}\in(0,1)$. When $\ell(\lambda)>n$, the cell
$(n+1,1)$ gives $p^n C_1(\lambda)\ge1$, and the inequality remains true.

For $0\le a\le1$, denote the numerator of
\eqref{eq:scaled-schur-sum} by $N_{n,a}(z)$ and its infinite-alphabet
version by $N_{\infty,a}(z)$. Equations
\eqref{eq:finite-alphabet-error} and \eqref{eq:content-moment-bound}
give, uniformly for $|z|\le R$,
\[
 |N_{n,a}(z)-N_{\infty,a}(z)|
 \le \frac{2\mathcal M\Phi_p(R)(1+R)}{(1-p)^2}\,p^{n+1}.
\]
Writing $\mathcal A_n(a)=\prod_{i,j<n}(1-ap^{i+j+1})^{-1}$, the
same estimate with $R=0$ gives
\[
 |\mathcal A_n(a)-\mathcal A_\infty(a)|
 \le \frac{2\mathcal M}{(1-p)^2}\,p^{n+1}.
\]
Since $\mathcal A_n(a)\ge1$, $|\Phi_p(z)|\le\Phi_p(R)$, and
$N_{\infty,a}=\mathcal A_\infty(a)\Phi_p$, division proves
\eqref{eq:scaled-limit-explicit} for $a=p^r$. Cauchy's derivative
estimates on a slightly larger disk prove the derivative assertion.
\end{proof}

\begin{corollary}[Zero clusters and the closest remaining zero]
\label{cor:zero-clusters}
Fix $k\ge1$, and choose a sufficiently small disk centered at $-p^{-k}$
that contains no other zero of $\Phi_p$. For all sufficiently large $n$,
independently of $r\ge0$, this disk contains exactly $k$ zeros of
$F_{n,r}$, counted with multiplicity. All these zeros tend to $-p^{-k}$,
uniformly in $r$.

In particular, for all sufficiently large $n$, each $\Psi_{n,r}$ has
a unique zero of smallest modulus, denoted $\beta_{n,r}$. It is simple,
real, and negative, and
\begin{equation}\label{eq:closest-zero-leading}
 \sup_{r\ge0}|p^{n+r}\beta_{n,r}+p^{-1}|\longrightarrow0.
\end{equation}
Moreover, $\Psi_{n,r}\to1$ locally uniformly in the unscaled $\alpha$
plane, uniformly in $r$. Hence, on every fixed bounded disk, the zeros
of $H_{n,r}$ are eventually exactly those in
\eqref{eq:forced-alpha-orders}, with the stated multiplicities.
\end{corollary}

\begin{proof}
The product in \eqref{eq:infinite-normalizer-and-Phi} has precisely the
zeros $-p^{-k}$, of multiplicity $k$. Apply Rouch\'e's theorem to the
boundary of each chosen disk, using \cref{thm:scaled-limit}.
The same argument on $|z|=\rho$, with $p^{-1}<\rho<p^{-2}$,
shows that there is exactly one zero in that disk. The zero near
$-p^{-1}$ is therefore the unique closest one. It is simple and real
because the functions have real coefficients and the disk is invariant
under complex conjugation. Its limit is negative, proving the sign
assertion and \eqref{eq:closest-zero-leading}.

Finally,
$\Psi_{n,r}(\alpha)=F_{n,r}(\alpha p^{n+r})$ and $\Phi_p(0)=1$.
The quantitative theorem therefore gives convergence to one uniformly
on fixed $\alpha$-disks and uniformly in $r$. Combining this with
\cref{thm:forced-alpha-orders} proves the last assertion.
\end{proof}

\begin{remark}
For $k>1$, the cluster statement alone does not assert that the $k$
zeros remain a multiple zero or that they are real. The splitting and
reality questions require additional estimates; the double cluster is
analyzed below.
\end{remark}

\subsection{The full unshifted leading constant and its first correction}

\begin{theorem}[Fixed-parameter refinement]\label{thm:fixed-alpha-refinement}
For each $R>0$, uniformly on $|\alpha|\le R$,
\begin{equation}\label{eq:fixed-alpha-refinement}
 \frac{H_{n,0}(\alpha)}
 {p^{K_{n,0}}P_\infty^{2n}\mathcal L_{n,0}(\alpha)}
 =\mathcal M^3\left(
  1+\frac{(\alpha-4)p^{n+1}}{(1-p)^2}
  +O_{p,R}(p^{2n})\right).
\end{equation}
At the forced zeros the quotient is understood by its entire extension.
In particular, the first correction vanishes at $\alpha=4$.
\end{theorem}

\begin{proof}
Write $F_n=F_{n,0}$. We have $F_n(0)=1$ exactly. By the derivative
part of \cref{thm:scaled-limit},
\[
 F_n'(0)=\Phi_p'(0)+O_p(p^n)
 =\frac{p}{(1-p)^2}+O_p(p^n),
\]
and $F_n''$ is uniformly bounded on each fixed disk. Taylor expansion
at zero therefore gives
\begin{equation}\label{eq:fixed-alpha-Psi-correction}
 \Psi_{n,0}(\alpha)=F_n(\alpha p^n)
 =1+\frac{\alpha p^{n+1}}{(1-p)^2}+O_{p,R}(p^{2n}).
\end{equation}
The exact Cauchy evaluation can also be written as
\[
 \mathcal C_{n,0}
 =p^{K_{n,0}}
   \frac{\prod_{j=1}^{n-1}\pi_j^3}
        {\prod_{j=n}^{2n-1}\pi_j}.
\]
Taking real logarithms of its positive factors yields the exact identity
\[
 \log\frac{\mathcal C_{n,0}}
 {p^{K_{n,0}}P_\infty^{2n}\mathcal M^3}
 =4\sum_{m\ge1}m\log(1-p^{n+m})
  -\sum_{m\ge1}m\log(1-p^{2n+m}).
\]
Consequently,
\[
 \frac{\mathcal C_{n,0}}
 {p^{K_{n,0}}P_\infty^{2n}\mathcal M^3}
 =1-\frac{4p^{n+1}}{(1-p)^2}+O_p(p^{2n}).
\]
Multiplying this estimate by \eqref{eq:fixed-alpha-Psi-correction}
proves \eqref{eq:fixed-alpha-refinement}.
\end{proof}

\begin{corollary}[Unnormalized complex-parameter asymptotic]
\label{cor:unnormalized-full-asymptotic}
Fix $\alpha\in\C\setminus
\bigl(\{0\}\cup\{-p^{-m}:m\ge1\}\bigr)$, and put
\[
 \mathcal R_p(\alpha)
 =\prod_{m\ge1}(1+\alpha p^m)^{-m}.
\]
Then
\begin{equation}\label{eq:unnormalized-full-asymptotic}
 \begin{split}
 H_{n,0}(\alpha)
 ={}&\alpha^{n^2}p^{K_{n,0}}
 \bigl(E_q(\alpha)P_\infty^2\bigr)^n
 \mathcal M^3\mathcal R_p(\alpha)\\
 &\times\left(1+
 \frac{2(\alpha-2)p^{n+1}}{(1-p)^2}
 +O_{p,\alpha}(p^{2n})\right).
 \end{split}
\end{equation}
\end{corollary}

\begin{proof}
The absolutely convergent products give the exact identity
\[
 \frac{\mathcal L_{n,0}(\alpha)}
 {\alpha^{n^2}E_q(\alpha)^n\mathcal R_p(\alpha)}
 =\prod_{s\ge1}(1+\alpha p^{n+s})^s
 =1+\frac{\alpha p^{n+1}}{(1-p)^2}
   +O_{p,\alpha}(p^{2n}).
\]
Multiplication by \eqref{eq:fixed-alpha-refinement} proves the claim.
\end{proof}

These are analytic refinements of the determinant estimates; they do
not alter the cubic arithmetic comparison or assert a smaller
irrationality threshold.


\section{A classical moment model and the splitting of a double limiting zero}
\label{mx:section}

The specialization $\mu=0$ admits a classical orthogonal-polynomial
description.  We use the $q$-Meixner formulas of
Koekoek and Swarttouw~\cite[\S3.13]{KS}; the identification below also
specifies all normalizations needed for the determinant calculations.
Throughout this section $q>1$, $p=q^{-1}$, and
\begin{equation}\label{mx:notation}
 (a;p)_j=\prod_{s=0}^{j-1}(1-ap^s),\qquad
 \pi_j=(p;p)_j,\qquad
 P_\infty=(p;p)_\infty,\qquad
 \mathcal M=\prod_{s=1}^{\infty}(1-p^s)^{-s}.
\end{equation}
All empty products equal $1$.

\subsection{The zero-value moment functional}

Define the linear functional $L_\alpha$ on polynomials by
\begin{equation}\label{mx:functional}
 L_\alpha[x^s]=-v_s(q,\alpha,0),\qquad s\ge0.
\end{equation}
Its moments are polynomials in $\alpha$ and satisfy
\begin{equation}\label{mx:moment-recurrence}
 m_0=1,\qquad m_s=(q^s-1)m_{s-1}+\alpha^s\quad(s\ge1).
\end{equation}

\begin{lemma}[The classical $q$-Meixner specialization]\label{mx:meixner}
For $\alpha>0$, the functional $L_\alpha$ has the positive discrete
representation
\begin{equation}\label{mx:positive-measure}
 L_\alpha[f]
 =\sum_{s=0}^{\infty}\lambda_s f(\alpha p^{-s}),\qquad
 \lambda_s=
 \frac{\alpha^{-s}p^{\binom{s}{2}}}
 {(-\alpha^{-1};p)_{s+1}}.
\end{equation}
Its monic orthogonal polynomials are
\begin{align}
 P_j^\alpha(x)
 &=(-1)^j\pi_jp^{-j^2}
 M_j(x/\alpha;1,\alpha^{-1};p)\label{mx:monic-polynomial}\\
 &=(-1)^j\pi_jp^{-j^2}
 \sum_{\ell=0}^{j}
 \frac{(p^{-j};p)_\ell(x/\alpha;p)_\ell}{\pi_\ell^2}
 (-\alpha p^{j+1})^\ell.\nonumber
\end{align}
Here $M_j$ is the $q$-Meixner polynomial in the normalization of
\cite[\S3.13]{KS}.  Their squared norms are
\begin{align}
 h_j(\alpha):=L_\alpha[(P_j^\alpha)^2]
 &=\pi_j^2p^{-j(2j+1)}(-\alpha p;p)_j\label{mx:norm}\\
 &=q^{\binom j2}
 \prod_{a=1}^{j}(q^a-1)^2
 \prod_{a=1}^{j}(\alpha+q^a).\nonumber
\end{align}
The polynomials in \eqref{mx:monic-polynomial} extend polynomially to
$\alpha=0$.  The identities
$L_\alpha[P_i^\alpha P_j^\alpha]=\delta_{ij}h_j(\alpha)$
hold for every complex $\alpha$, with possible degeneracy when a norm
vanishes.
\end{lemma}

\begin{proof}
Put
$t_s=\alpha^{-s}p^{\binom{s}{2}}/(-\alpha^{-1};p)_s$.
Then $\lambda_s=t_s-t_{s+1}$, $t_0=1$, and $t_s\to0$.
Thus $\sum_s\lambda_s=1$.
All polynomial moments in \eqref{mx:positive-measure} converge absolutely.
Moreover,
\[
 (\alpha+1)\lambda_0=\alpha,\qquad
 (1+\alpha p^{-s})\lambda_s=q\lambda_{s-1}\quad(s\ge1).
\]
Multiplication by $(\alpha p^{-s})^{k-1}$ and summation show that these
moments satisfy \eqref{mx:moment-recurrence}.  This proves
\eqref{mx:positive-measure}.

The standard $q$-Meixner orthogonality formula, specialized to base $p$,
$b=1$, and $c=\alpha^{-1}$, reads in this probability normalization
\[
 \sum_{s=0}^{\infty}\lambda_s
 M_i(p^{-s};1,\alpha^{-1};p)
 M_j(p^{-s};1,\alpha^{-1};p)
 =\delta_{ij}p^{-j}(-\alpha p;p)_j.
\]
The term $\ell=j$ in the terminating series shows that its leading
coefficient as a polynomial in $x$ is
$(-1)^jp^{j^2}/\pi_j$.  The multiplier in
\eqref{mx:monic-polynomial} therefore makes it monic, and multiplication
of the displayed norm by $\pi_j^2p^{-2j^2}$ gives \eqref{mx:norm}.
Finally, every expression
$\alpha^\ell(x/\alpha;p)_\ell$
is a polynomial in $(x,\alpha)$.  Both sides of each orthogonality
identity are consequently polynomials in $\alpha$.  Their validity for
$\alpha>0$ extends them to all complex $\alpha$.
\end{proof}

\begin{proposition}[Exact zero-value determinant]\label{mx:zero-value-product}
For every $n\ge1$,
\begin{equation}\label{mx:zero-value-factorization}
 V_n(q,\alpha,0)
 =(-1)^nq^{\binom n3}
 \prod_{d=1}^{n-1}(q^d-1)^{2(n-d)}
 \prod_{m=1}^{n-1}(\alpha+q^m)^{n-m}.
\end{equation}
In particular, its degree in $\alpha$ is $\binom n2$.
\end{proposition}

\begin{proof}
Changing from monomials to monic orthogonal polynomials is a triangular
change of basis with determinant $1$.  Thus the Hankel determinant of
$L_\alpha$ is $\prod_{j=0}^{n-1}h_j(\alpha)$.
The moments of $V_n(q,\alpha,0)$ are the negatives of those of
$L_\alpha$, so its determinant is
$(-1)^n\prod_{j=0}^{n-1}h_j(\alpha)$.
Substitution of \eqref{mx:norm}, followed by collecting each factor,
gives \eqref{mx:zero-value-factorization} for $\alpha>0$.
The polynomial identity extends to every $\alpha$.
\end{proof}

We shall also use the fixed positive functional
\begin{equation}\label{mx:B-functional}
 B=L_0,\qquad
 B[x^s]=\prod_{a=1}^{s}(q^a-1)
       =p^{-s(s+1)/2}\pi_s.
\end{equation}
It is positive definite on real polynomials: at $\alpha=0$ the
polynomials $P_j^0$ remain monic and orthogonal, and their norms
$\pi_j^2p^{-j(2j+1)}$ are strictly positive.

\subsection{Logarithmic derivatives at the product zeros}

The following estimates concern the derivative in the independent
variable $\mu$, before the specialization $\mu=E_q(\alpha)$.
We write a logarithmic derivative as a quotient of derivatives, so
that no choice of a complex logarithm is needed.

\begin{theorem}\label{mx:log-derivative}
Fix an integer $k\ge1$.  For $n\ge1$ set $N=n+k$ and
$\alpha_N=-p^{-N}$.  Then $V_n(q,\alpha_N,0)\ne0$, and
\begin{equation}\label{mx:log-derivative-asymptotic}
 T_{n,k}:=
 \frac{\partial_\mu V_n(q,\alpha_N,0)}{V_n(q,\alpha_N,0)}
 \sim
 \frac{(-1)^n}{P_\infty\pi_{k-1}}
 p^{-\{n^2+(2k-3)n-2k+2\}/2}.
\end{equation}
\end{theorem}

\begin{proof}
Solving the homogeneous part of the recurrence for $v_s$ gives
\[
 v_s(q,\alpha,\mu)=-L_\alpha[x^s]+\mu B[x^s].
\]
For fixed $\alpha=\alpha_N$, pass in this moment matrix to the monic
basis $P_0^\alpha,\ldots,P_{n-1}^\alpha$.
At $\mu=0$ the matrix is diagonal with entries $-h_j(\alpha)$.
These entries are nonzero because $j<n<N$.
Differentiation of the determinant therefore gives the exact identity
\begin{equation}\label{mx:trace-formula}
 T_{n,k}=-\sum_{j=0}^{n-1}
 \frac{B[(P_j^{\alpha_N})^2]}{h_j(\alpha_N)}.
\end{equation}

Put $r=N-j$ and
\[
 A_j=(-1)^j\alpha_N^jp^{-\binom j2},\qquad
 R_{j,N}(x)=P_j^{\alpha_N}(x)/A_j.
\]
Reversing the terminating sum in \eqref{mx:monic-polynomial}, using
$(p^{-j};p)_s=(-1)^sp^{-js+\binom s2}\pi_j/\pi_{j-s}$,
gives the exact formula
\begin{equation}\label{mx:reversed-polynomial}
 R_{j,N}(x)=
 \sum_{\ell=0}^{j}
 \frac{(-1)^\ell p^{r\ell+\binom\ell2}}{\pi_\ell}
 \left(\frac{\pi_j}{\pi_{j-\ell}}\right)^2
 (-xp^N;p)_{j-\ell}.
\end{equation}
We next prove both a limit and a bound for its $B$-norm.

For $0\le s\le N-2$, write $U_s(x)=(-xp^N;p)_s$.
The finite $p$-binomial formula gives
\[
 [x^u]U_s(x)
 =p^{Nu+\binom u2}\frac{\pi_s}{\pi_u\pi_{s-u}}
 \le P_\infty^{-1}p^{Nu+\binom u2}
 \quad(0\le u\le s).
\]
All its coefficients are nonnegative.  In the product of its
$x^u$ and $x^v$ terms, the moment bound from
\eqref{mx:B-functional} gives the exponent
\[
 N(u+v)+\binom u2+\binom v2
 -\frac{(u+v)(u+v+1)}2
 =(N-1)(u+v)-uv.
\]
Since $u,v\le N-2$,
$uv\le (N-2)(u+v)/2$, and this exponent is at least
$N(u+v)/2$.  Removing the constant term, whose value is exactly $1$,
we obtain the uniform bound
\begin{equation}\label{mx:U-bound}
 0\le B[U_s^2]-1
 \le P_\infty^{-2}
 \left\{(1-p^{N/2})^{-2}-1\right\}
 \longrightarrow0.
\end{equation}
Also $B[(U_s-1)^2]\le B[U_s^2]-1$, because
$B[U_s]\ge1$.

In \eqref{mx:reversed-polynomial} we have
$r\ge k+1\ge2$ and $j\le N-2$.
The absolute sum of its scalar coefficients is bounded independently
of $N,j,r$ by
\[
 C_p=\sum_{\ell=0}^{\infty}
 \frac{p^{2\ell+\binom\ell2}}{\pi_\ell}<\infty.
\]
Let $S_{j,r}$ denote their signed sum, obtained by replacing each
$U_{j-\ell}$ by $1$.
The triangle inequality for the norm
$\|f\|_B=B[f^2]^{1/2}$ and \eqref{mx:U-bound} imply
\[
 \|R_{j,N}-S_{j,r}\|_B\longrightarrow0
 \quad\hbox{uniformly for }0\le j\le N-2.
\]
Consequently $B[R_{j,N}^2]$ is uniformly bounded.
For fixed $r\ge2$, as $N\to\infty$ and $j=N-r$,
dominated convergence in the scalar sum and Euler's product identity
give
\[
 S_{j,r}\longrightarrow
 \sum_{\ell=0}^{\infty}
 \frac{(-1)^\ell p^{r\ell+\binom\ell2}}{\pi_\ell}
 =(p^r;p)_\infty.
\]
It follows that
\begin{equation}\label{mx:B-norm-limit}
 B[R_{N-r,N}^2]\longrightarrow(p^r;p)_\infty^2
 \qquad\hbox{for each fixed }r\ge2.
\end{equation}

An exact simplification of \eqref{mx:norm} yields
\begin{equation}\label{mx:norm-ratio}
 \frac{A_j^2}{h_j(\alpha_N)}
 =\frac{(-1)^j
 p^{-(N^2-3N)/2+r(r-3)/2}}
 {\pi_j^2(p^r;p)_j}.
\end{equation}
Substitution into \eqref{mx:trace-formula} shows that
\begin{equation}\label{mx:scaled-trace}
 \frac{T_{n,k}}{(-1)^{N+1}p^{-(N^2-3N)/2}}
 =\sum_{r=k+1}^{N}
 (-1)^rp^{r(r-3)/2}
 \frac{B[R_{N-r,N}^2]}
 {\pi_{N-r}^2(p^r;p)_{N-r}}.
\end{equation}
Each denominator is at least $P_\infty^3$, and the numerators just
proved are uniformly bounded.  The summable sequence
$C p^{r(r-3)/2}$ therefore dominates the summands.
Using \eqref{mx:B-norm-limit}, the right side tends to
\begin{align}
 \Sigma_k
 &=\frac1{P_\infty^2}
 \sum_{r=k+1}^{\infty}
 (-1)^rp^{r(r-3)/2}(p^r;p)_\infty\label{mx:Euler-tail}\\
 &=-\frac{p^{-1}}{P_\infty}
 \sum_{\ell=k}^{\infty}
 \frac{(-1)^\ell p^{\binom\ell2}}{\pi_\ell}.
 \nonumber
\end{align}
For clarity, the cancellation in this last sum is explicit:
Euler's identity at $1$ gives the full sum equal to $0$, while
\[
 \sum_{\ell=0}^{k-1}
 \frac{(-1)^\ell p^{\binom\ell2}}{\pi_\ell}
 =\frac{(-1)^{k-1}p^{\binom k2}}{\pi_{k-1}}.
\]
The finite identity follows by subtracting its instances for
successive $k$.  Hence
\[
 \Sigma_k=
 \frac{(-1)^{k-1}p^{\binom k2-1}}
 {P_\infty\pi_{k-1}}\ne0.
\]
Combining this value with \eqref{mx:scaled-trace} and $N=n+k$
gives \eqref{mx:log-derivative-asymptotic}.
\end{proof}

\subsection{Exact center values and their slopes}

We use the normalization introduced earlier,
\[
 F_n(z)=F_{n,0}(z)
 =\Psi_{n,0}(p^{-n}z),\qquad
 \Phi_p(z)=\prod_{\ell=1}^{\infty}(1+zp^\ell)^\ell.
\]
In particular, \cref{thm:scaled-limit} gives local uniform convergence
$F_n\to\Phi_p$, and hence local uniform convergence of every fixed
derivative.

\begin{corollary}[Values and slopes at the limiting zeros]
\label{mx:center-slopes}
For every integer $k\ge1$ and every $n\ge1$,
\begin{equation}\label{mx:exact-center}
 F_n(-p^{-k})=\frac{p^{kn^2}}{\mathcal A_{n,0}}.
\end{equation}
For fixed $k$, as $n\to\infty$,
\begin{equation}\label{mx:center-slope}
 F_n'(-p^{-k})
 \sim
 \frac{(-1)^{k-1}P_\infty}{\mathcal M\pi_{k-1}}
 p^{(k-1)n(n-2)+(-k^2+5k-2)/2}.
\end{equation}
\end{corollary}

\begin{proof}
The Cauchy determinant in the unshifted normalization is
\[
 \mathcal C_{n,0}
 =p^{n(n+1)(2n+1)/6}
 \prod_{d=1}^{n-1}(1-p^d)^{2(n-d)}\mathcal A_{n,0}.
\]
Writing \eqref{mx:zero-value-factorization} in powers of $p$ and
cancelling the common factors gives, with $\alpha=p^{-n}z$,
\begin{equation}\label{mx:normalized-zero-value}
 \frac{V_n(q,p^{-n}z,0)}
 {\mathcal L_{n,0}(p^{-n}z)\mathcal C_{n,0}}
 =\frac{(-1)^nz^{-n^2}}{\mathcal A_{n,0}}.
\end{equation}
The identity is first read away from the vanishing factors and
then by cancellation.  At $z_k=-p^{-k}$ we have
$E_q(p^{-n}z_k)=0$, and \eqref{mx:exact-center} follows.

Set $\mu_n(z)=E_q(p^{-n}z)$ and $N=n+k$.
Differentiation of its product at the simple zero $z_k$ gives
\begin{equation}\label{mx:E-derivative}
 \mu_n'(z_k)
 =(-1)^{N-1}p^{k-N(N-1)/2}\pi_{N-1}P_\infty.
\end{equation}
The chain rule applied to \eqref{mx:normalized-zero-value} gives the
exact relation
\begin{equation}\label{mx:chain-rule}
 \frac{F_n'(z_k)}{F_n(z_k)}
 =\mu_n'(z_k)T_{n,k}-\frac{n^2}{z_k}.
\end{equation}
Now $\mathcal A_{n,0}\to\mathcal M$ and
$\pi_{N-1}\to P_\infty$.
Insert \eqref{mx:log-derivative-asymptotic} and
\eqref{mx:E-derivative} into \eqref{mx:chain-rule}.
The term $n^2/z_k$ is negligible compared with
$\mu_n'(z_k)T_{n,k}$.  Collecting its sign and exponent gives
\eqref{mx:center-slope}.
\end{proof}

\subsection{Two different displacement scales at the second center}

The limiting zero $-p^{-2}$ has multiplicity two in $\Phi_p$.
Its two approximating zeros do not have the same displacement scale.

\begin{theorem}[Asymmetric splitting at the second center]
\label{mx:double-splitting}
For all sufficiently large $n$, $F_n$ has exactly two zeros in a
sufficiently small fixed disk about $-p^{-2}$.
They are real and simple, and can be labelled so that
\begin{align}
 z_{n,+}
 &=-p^{-2}
 +\frac{1-p}{p^2P_\infty}
 p^{n^2+2n}(1+o(1)),\label{mx:near-root}\\
 z_{n,-}
 &=-p^{-2}
 -\frac{1}{p(1-p)^2P_\infty}
 p^{n^2-2n}(1+o(1)).\label{mx:far-root}
\end{align}
Equivalently, the corresponding zeros of $H_{n,0}$ are
\begin{align}
 \alpha_{n,+}
 &=-q^{n+2}
 +\frac{1-p}{P_\infty}
 p^{n^2+n-2}(1+o(1)),\label{mx:near-alpha}\\
 \alpha_{n,-}
 &=-q^{n+2}
 -\frac{1}{(1-p)^2P_\infty}
 p^{n^2-3n-1}(1+o(1)).\label{mx:far-alpha}
\end{align}
\end{theorem}

\begin{proof}
Put $z_0=-p^{-2}$.  The product defining $\Phi_p$ gives
\begin{equation}\label{mx:second-limit-derivative}
 \frac{\Phi_p''(z_0)}2
 =p^4(1-p^{-1})
 \prod_{\ell=3}^{\infty}(1-p^{\ell-2})^\ell
 =-\frac{p^3(1-p)P_\infty^2}{\mathcal M}.
\end{equation}
Write
\[
 a=\mathcal M^{-1},\qquad
 b=\frac{p^2P_\infty}{(1-p)\mathcal M},\qquad
 c=\frac{p^3(1-p)P_\infty^2}{\mathcal M};
\]
all three constants are positive.
The exact center value, the slope asymptotic, and derivative
convergence give
\begin{align*}
 F_n(z_0)&=a p^{2n^2}(1+o(1)),\\
 F_n'(z_0)&=-b p^{n^2-2n}(1+o(1)),\\
 \frac{F_n''(z_0)}2&=-c+o(1).
\end{align*}
On a fixed disk about $z_0$, the third derivatives are uniformly
bounded by local uniform convergence and Cauchy's estimates.
Thus, uniformly for sufficiently small $\delta$,
\begin{equation}\label{mx:local-Taylor}
 F_n(z_0+\delta)
 =F_n(z_0)+F_n'(z_0)\delta
 +\tfrac12F_n''(z_0)\delta^2+O(\delta^3).
\end{equation}

First let $\varepsilon_n=p^{n^2-2n}$ and set
$\delta=\varepsilon_nt$.
Division of \eqref{mx:local-Taylor} by $\varepsilon_n^2$ gives
local uniform convergence in $t$ to
\[
 -bt-ct^2.
\]
Indeed, the constant term after division is
$a p^{4n}(1+o(1))$, and the third-order remainder is
$O(\varepsilon_n)$ on bounded $t$-sets.
The limiting zero $-b/c$ is simple.
Rouch\'e's theorem on a small circle about that zero gives exactly
one zero with
\[
 \delta=-\frac bc\varepsilon_n(1+o(1)).
\]
This is \eqref{mx:far-root}.

To resolve the zero that is $o(\varepsilon_n)$ from the center,
let $\eta_n=p^{n^2+2n}$ and set $\delta=\eta_nu$.
Division of \eqref{mx:local-Taylor} by $p^{2n^2}$ now gives
local uniform convergence in $u$ to
\[
 a-bu.
\]
The quadratic term is $O(p^{4n})$ on bounded $u$-sets and the
third-order term tends to zero.
Rouch\'e's theorem about the simple zero $a/b$ gives exactly
one zero with
\[
 \delta=\frac ab\eta_n(1+o(1)),
\]
which is \eqref{mx:near-root}.

Choose a fixed disk about $z_0$ containing no other zero of
$\Phi_p$.  By \cref{thm:scaled-limit} and Rouch\'e's theorem,
$F_n$ has precisely two zeros there, counted with multiplicity,
for all sufficiently large $n$.
The two zeros just constructed exhaust them and are simple.
The coefficients of $F_n$ are real, so conjugation preserves its
zeros.  Each of the two small disks used above is invariant under
conjugation and contains exactly one zero; hence both zeros are real.
The factors divided out in the definition of $F_n$ do not vanish
near this center.  Finally, multiplication by $p^{-n}$ gives
\eqref{mx:near-alpha}--\eqref{mx:far-alpha}.
\end{proof}

\begin{remark}\label{mx:scope}
The fixed-$k$ slope formula does not by itself determine the full
splitting of a limiting zero of multiplicity $k\ge3$.
Such a conclusion would require suitably scaled estimates for
higher derivatives.  The present section proves the center values
and slopes for every fixed $k$, and the complete two-zero splitting
only for $k=2$.
\end{remark}


\section{Dimension reduction and the first residual zero for fixed shifts}
\label{sz:section}

Throughout this section, $q>1$, $p=q^{-1}$, and
$\pi_j=(p;p)_j$, with $\pi_0=1$.  We retain the residual functions
$\Psi_{n,r}$ and their rescalings
\[
 F_{n,r}(z)=\Psi_{n,r}(p^{-n-r}z)
\]
from \cref{thm:scaled-limit}.  The shift $r$ is fixed when $n$ tends to
infinity.  We give an exact reduction of a special value of the original
$n\times n$ determinant to an $r\times r$ determinant, and use it to locate
the first zero of $\Psi_{n,r}$ beyond the explicitly removed factors.

\begin{lemma}[A terminating sum and its limit]\label{sz:sum-limit}
For $N\ge0$ and an integer $a$, define
\begin{equation}\label{sz:def-S}
 S_N(a)=\sum_{\ell=0}^{N}
 \frac{(-1)^\ell p^{\binom{\ell}{2}+a\ell}\pi_N^2}
 {\pi_\ell\pi_{N-\ell}^2}.
\end{equation}
For every fixed integer $a$, including $a\le0$,
\begin{equation}\label{sz:S-limit}
 S_N(a)=(p^a;p)_\infty+O_{p,a}(p^N).
\end{equation}
In particular, the limiting value in \eqref{sz:S-limit} is zero if $a\le0$.
\end{lemma}

\begin{proof}
For $0\le\ell\le N$, put
\[
 R_{N,\ell}=\frac{\pi_N^2}{\pi_{N-\ell}^2}
 =\prod_{h=N-\ell+1}^{N}(1-p^h)^2.
\]
Since all factors lie in $[0,1]$,
\begin{equation}\label{sz:ratio-estimate}
 0\le 1-R_{N,\ell}
 \le 2\sum_{h=N-\ell+1}^{N}p^h
 \le\frac{2p^{N-\ell+1}}{1-p}.
\end{equation}
Euler's identity, valid for every complex argument, gives
\[
 (p^a;p)_\infty
 =\sum_{\ell=0}^{\infty}
 \frac{(-1)^\ell p^{\binom{\ell}{2}+a\ell}}{\pi_\ell}.
\]
The difference contributed by $0\le\ell\le N$ is at most
\[
 \frac{2p^{N+1}}{1-p}
 \sum_{\ell=0}^{\infty}
 \frac{p^{\binom{\ell}{2}+(a-1)\ell}}{\pi_\ell}
 =O_{p,a}(p^N).
\]
This series converges for every fixed $a$, since
$\pi_\ell\ge P_\infty>0$ and its numerator has quadratic decay in
$\ell$.  The omitted tail of Euler's series is also $O_{p,a}(p^N)$,
by the same quadratic decay.  This proves the assertion, without a
positivity assumption on $a$.
\end{proof}

We first recall the determinant identity needed for the reduction and
include an algebraic proof.  For a linear functional $\ell$ on
polynomials, write
\[
 \Delta_n^{(r)}=\det\bigl(\ell[x^{i+j+r}]\bigr)_{0\le i,j<n}.
\]

\begin{lemma}[Confluent Christoffel identity]\label{sz:christoffel}
Suppose that the moment functional has monic orthogonal polynomials
$Q_0,Q_1,\ldots,Q_{n+r-1}$ and that the required unshifted moment
determinants are nonzero.  Then
\begin{equation}\label{sz:christoffel-formula}
 \frac{\Delta_n^{(r)}}{\Delta_n^{(0)}}
 =(-1)^{nr}
 \det\bigl([x^i]Q_{n+j}(x)\bigr)_{0\le i,j<r}.
\end{equation}
For $r=0$, the determinant on the right is interpreted as $1$.
\end{lemma}

\begin{proof}
Let $\mathcal V_s(u)=\prod_{1\le i<j\le s}(u_j-u_i)$.
Apply the tensor product functional $\ell^{\otimes n}$ to the variables
$x_1,\ldots,x_n$.  For distinct auxiliary variables $t_1,\ldots,t_r$,
the Vandermonde factorization and orthogonality give
\begin{equation}\label{sz:christoffel-distinct}
 \frac{\ell^{\otimes n}\!\left[
 \mathcal V_n(x)^2\prod_{a=1}^{r}\prod_{i=1}^{n}(t_a-x_i)
 \right]}{n!\Delta_n^{(0)}}
 =\frac{\det\bigl(Q_{n+j}(t_a)\bigr)_{
 1\le a\le r,\,0\le j<r}}{\mathcal V_r(t)}.
\end{equation}
Here is a direct justification.  Factor
\[
 \mathcal V_{n+r}(x_1,\ldots,x_n,t_1,\ldots,t_r)
 =\mathcal V_n(x)\mathcal V_r(t)
 \prod_{a=1}^{r}\prod_{i=1}^{n}(t_a-x_i).
\]
In the determinant defining the left side, replace the monomials by the
monic polynomials $Q_0,\ldots,Q_{n+r-1}$.  Expand along the $n$ rows
containing the $x_i$.  After applying $\ell^{\otimes n}$ against
$\mathcal V_n(x)$, orthogonality eliminates every selection of columns
except $0,\ldots,n-1$.  The surviving factor is
$n!\Delta_n^{(0)}$, which proves \eqref{sz:christoffel-distinct}.

Let all $t_a$ tend to zero.  The quotient on the right tends to the
coefficient determinant in \eqref{sz:christoffel-formula}; this follows
by expanding each $Q_{n+j}(t_a)$ at zero and taking the lowest alternating
term.  On the left, the inserted product becomes
$(-1)^{nr}\prod_i x_i^r$.  Expanding the two Vandermonde determinants
shows that its image under $\ell^{\otimes n}$ is
$(-1)^{nr}n!\Delta_n^{(r)}$.  The result follows.
\end{proof}

\begin{theorem}[Exact reduction at the first limiting zero]
\label{sz:reduction}
For integers $n\ge1$ and $r\ge0$, put
\[
 R_r=\frac{r(r-1)}2,\qquad T_r=\frac{r(r-1)(r-2)}6,
\]
and define the $r\times r$ matrix
\begin{equation}\label{sz:def-B}
 B_{n,r}=\bigl(p^{ij}S_{n+j}(r+1-i-j)\bigr)_{0\le i,j<r}.
\end{equation}
Then
\begin{equation}\label{sz:exact-reduction}
 \boxed{
 \mathcal A_{n,r}F_{n,r}(-q)
 =\frac{(-1)^{R_r}p^{n^2-T_r}}
 {\prod_{i=0}^{r-1}\pi_i}\det B_{n,r}.}
\end{equation}
Equivalently, the residual function in \eqref{sz:exact-reduction} is
evaluated at $\alpha=-q^{n+r+1}$.  Empty determinants and products give
the exact identity
\begin{equation}\label{sz:unshifted-special-value}
 \mathcal A_{n,0}F_{n,0}(-q)=p^{n^2}.
\end{equation}
\end{theorem}

\begin{proof}
Write $W_n^{(r)}(\alpha)=V_n^{(r)}(q,\alpha,0)$ and set
\[
 m=n+r+1,\qquad \alpha=-p^{-m}.
\]
At this value of $\alpha$, the product formula for $E_q$ gives
$E_q(\alpha)=0$, so the actual determinant agrees with
$W_n^{(r)}(\alpha)$.  The monic polynomials for the moment functional
with moments $-v_k(q,\alpha,0)$ are denoted by $P_N^\alpha$ as in
\cref{mx:meixner}.  Multiplication of all moments by $-1$ has no effect
on the quotient of two determinants of the same size.  Thus
\cref{sz:christoffel} gives
\begin{equation}\label{sz:W-ratio}
 \frac{W_n^{(r)}(\alpha)}{W_n^{(0)}(\alpha)}
 =(-1)^{nr}
 \det\bigl([x^i]P_{n+j}^\alpha(x)\bigr)_{0\le i,j<r}.
\end{equation}
The identity follows first for positive $\alpha$, where the moment
functional is positive, and then as an identity of rational functions
of $\alpha$.  At the present specialization the needed unshifted
determinants are nonzero: their factors $1+\alpha p^h$ have $h<m$,
by \cref{mx:zero-value-product}.

For $N=n+j$, the terminating formula in \cref{mx:meixner}, with the
summation index reversed, becomes
\begin{align}
 P_N^\alpha(x)
 &=a_N\sum_{\ell=0}^{N}
 \frac{(-1)^\ell p^{\binom{\ell}{2}+(m-N)\ell}\pi_N^2}
 {\pi_\ell\pi_{N-\ell}^2}
 (-xp^m;p)_{N-\ell},\label{sz:reversed-polynomial}\\
 a_N&=p^{-N(N-1)/2-mN}.\notag
\end{align}
Divide all the polynomials in the coefficient determinant by the
common power series
$G(x)=(-xp^m;p)_\infty$, whose constant term is $1$.
For the first $r$ coefficients this division is multiplication by a
lower triangular matrix with diagonal entries $1$; hence it leaves
their determinant unchanged.  Euler's reciprocal-product expansion
gives
\[
 \frac{(-xp^m;p)_{N-\ell}}{(-xp^m;p)_\infty}
 =\sum_{i=0}^{\infty}
 \frac{(-1)^i x^i p^{i(m+N-\ell)}}{\pi_i}.
\]
Consequently,
\begin{equation}\label{sz:polynomial-coefficient}
 [x^i]\frac{P_N^\alpha(x)}{G(x)}
 =a_N\frac{(-1)^ip^{i(m+N)}}{\pi_i}
 S_N(m-N-i).
\end{equation}
Extracting the row and column factors in
\eqref{sz:polynomial-coefficient} yields
\begin{equation}\label{sz:coefficient-prefactor}
 \det\bigl([x^i]P_{n+j}^\alpha(x)\bigr)_{i,j<r}
 =\frac{(-1)^{R_r}
 p^{-rn(n-1)/2-mrn-T_r}}{\prod_{i=0}^{r-1}\pi_i}
 \det B_{n,r}.
\end{equation}
For example, the exponent here is obtained from
\[
 \sum_{j=0}^{r-1}\left(-\frac{(n+j)(n+j-1)}2-m(n+j)\right)
 +(m+n)\sum_{i=0}^{r-1}i
 =-\frac{rn(n-1)}2-mrn-T_r.
\]

For completeness, we track the remaining normalization.  Put
\[
 d_n=\frac{n(n-1)(4n+1)}6,\quad
 J_n=\prod_{h=1}^{n-1}(1-p^h)^{2(n-h)},\quad
 L_n(\alpha)=\prod_{h=1}^{n-1}(1+\alpha p^h)^{n-h}.
\]
The zero-value product and the Cauchy determinant evaluation read
\[
 W_n^{(0)}(\alpha)=(-1)^np^{-d_n}J_nL_n(\alpha),
 \qquad
 \mathcal C_{n,r}=p^{K_{n,r}}J_n\mathcal A_{n,r}.
\]
Since $d_n+K_{n,r}=n^3+rn(n+1)/2$ and
$\mathcal L_{n,r}(\alpha)=\alpha^{n(n+r)}L_n(\alpha)$,
substitution of \eqref{sz:W-ratio} into the definition of $\Psi_{n,r}$
gives
\begin{equation}\label{sz:normalization-prefactor}
 \mathcal A_{n,r}F_{n,r}(-q)
 =p^{-n^3-rn(n+1)/2+mn(n+r)}
 \det\bigl([x^i]P_{n+j}^\alpha(x)\bigr)_{i,j<r}.
\end{equation}
The signs cancel because
$n+n(n+r)+nr=n(n+1)+2nr$ is even.
Combining \eqref{sz:coefficient-prefactor} and
\eqref{sz:normalization-prefactor}, the remaining exponent is
\[
 -n^3-\frac{rn(n+1)}2+mn(n+r)
 -\frac{rn(n-1)}2-mrn-T_r=n^2-T_r,
\]
where $m=n+r+1$.  This proves \eqref{sz:exact-reduction}.
\end{proof}

\begin{corollary}[Asymptotic value at the center]\label{sz:center-value}
For every fixed integer $r\ge0$, let
\begin{equation}\label{sz:def-center-constant}
 C_r=\prod_{a=2}^{r+1}(p^a;p)_\infty
 =\frac{P_\infty^r}{\prod_{j=1}^{r}\pi_j},\qquad C_0=1.
\end{equation}
Then
\begin{equation}\label{sz:center-asymptotic}
 \mathcal A_{n,r}F_{n,r}(-q)
 =p^{n^2}C_r\bigl(1+O_{p,r}(p^n)\bigr).
\end{equation}
\end{corollary}

\begin{proof}
By \cref{sz:sum-limit}, the entries of $B_{n,r}$ converge at rate
$O_{p,r}(p^n)$ to those of the matrix
\[
 B_{\infty,r}
 =\bigl(p^{ij}(p^{r+1-i-j};p)_\infty\bigr)_{0\le i,j<r}.
\]
For every entry, including those with a vanishing infinite product,
\begin{equation}\label{sz:newton-factorization}
 p^{ij}(p^{r+1-i-j};p)_\infty
 =(p^{r+1-j};p)_\infty
 \prod_{h=1}^{i}(p^j-p^{r+1-h}).
\end{equation}
After the first factor is extracted from each column, row $i$ consists
of the values at $1,p,\ldots,p^{r-1}$ of a monic polynomial of degree
$i$.  Its determinant is therefore the Vandermonde determinant.  Thus
\begin{align}
 \det B_{\infty,r}
 &=C_r\prod_{0\le j<k<r}(p^k-p^j)\notag\\
 &=(-1)^{R_r}p^{T_r}
 C_r\prod_{i=0}^{r-1}\pi_i.\label{sz:limit-determinant}
\end{align}
In particular this determinant is nonzero.  Since $r$ is fixed,
\eqref{sz:sum-limit} implies
$\det B_{n,r}=\det B_{\infty,r}(1+O_{p,r}(p^n))$.
The conclusion now follows from \cref{sz:reduction}.
\end{proof}

\begin{theorem}[First residual zero for each fixed shift]
\label{sz:first-zero}
Fix $q>1$ and an integer $r\ge0$.  For all sufficiently large $n$,
$\Psi_{n,r}$ has a unique zero of smallest modulus.  It is simple,
real, and negative.  Denote it by $\beta_{n,r}$.  Then
\begin{equation}\label{sz:first-zero-asymptotic}
 \boxed{
 \beta_{n,r}
 =-q^{n+r+1}
 -\frac{q^{-n^2+n+r+1}}{P_\infty\pi_r}
 \bigl(1+O_{q,r}(q^{-n})\bigr).}
\end{equation}
In particular, $\beta_{n,r}<-q^{n+r+1}$ for all sufficiently large
$n$.  The estimate concerns each fixed $r$; its implicit constant is
allowed to depend on $r$.
\end{theorem}

\begin{proof}
The limiting function in \cref{thm:scaled-limit} is
\[
 \Phi_p(z)=\prod_{\ell=1}^{\infty}(1+zp^\ell)^\ell.
\]
Its first zero is the simple zero $-q$, and the next zero is $-q^2$.
Choose $c$ with $q<c<q^2$.  Uniform convergence on $|z|=c$ and
Rouch\'e's theorem show that $F_{n,r}$ has exactly one zero, counted
with multiplicity, in $|z|<c$ for all sufficiently large $n$.
This zero, denoted $z_{n,r}$, tends to $-q$.  It is simple, and it is
real because $F_{n,r}$ takes real values on the real axis and its
nonreal zeros occur in conjugate pairs.  It is consequently the
unique zero of smallest modulus.

The products defining $\mathcal A_{n,r}$ and $\mathcal A_{\infty,r}$
give, for fixed $r$,
\begin{equation}\label{sz:A-tail}
 \mathcal A_{n,r}=\mathcal A_{\infty,r}
 \bigl(1+O_{p,r}(p^n)\bigr).
\end{equation}
Indeed, the sum of $p^{r+i+j+1}$ over the omitted pairs with
$i\ge n$ or $j\ge n$ is $O_{p,r}(p^n)$, which also bounds the
logarithmic product tail.  Therefore \cref{sz:center-value} gives
\begin{equation}\label{sz:F-center}
 F_{n,r}(-q)=\frac{C_r}{\mathcal A_{\infty,r}}
 p^{n^2}\bigl(1+O_{p,r}(p^n)\bigr)>0
\end{equation}
for all sufficiently large $n$.

Differentiating the limiting product at its simple zero yields
\begin{equation}\label{sz:Phi-derivative}
 \Phi_p'(-q)=\frac{pP_\infty}{\mathcal M}>0.
\end{equation}
By the derivative form of \cref{thm:scaled-limit}, $F_{n,r}'$ differs
from $\Phi_p'$ by $O_{p,r}(p^n)$ on a fixed neighborhood of $-q$.
The real mean value theorem first gives
$z_{n,r}+q=O_{p,r}(p^{n^2})$, since the derivative there is bounded
away from zero.  Applied once more with the derivative estimate, it
gives
\begin{align}
 z_{n,r}+q
 &=-\frac{\mathcal M C_r}{P_\infty\mathcal A_{\infty,r}}
 p^{n^2-1}\bigl(1+O_{p,r}(p^n)\bigr)\notag\\
 &=-\frac{p^{n^2-1}}{P_\infty\pi_r}
 \bigl(1+O_{p,r}(p^n)\bigr).\label{sz:scaled-displacement}
\end{align}
For the second equality, cancellation of the convergent products gives
\begin{equation}\label{sz:constant-cancellation}
 \frac{\mathcal M C_r}{\mathcal A_{\infty,r}}=\frac1{\pi_r}.
\end{equation}
Finally $\beta_{n,r}=p^{-n-r}z_{n,r}$, which proves
\eqref{sz:first-zero-asymptotic}, including the sign of the displacement.
\end{proof}


\section{Conclusions and remaining questions}\label{sec:conclusions}

The formal moment representation determines the generic polynomial order,
and the positive geometric representation determines the actual cubic
decay. With the prescribed cyclotomic divisor, these yield an explicit
restriction on proportional shifts: the arithmetic comparison cannot
improve its unshifted critical values $12/7$ and $6$.

The analytic normalization retains information beyond these cubic rates.
Its Schur expansion gives a uniform entire scaling limit and a fixed-parameter
asymptotic with its first correction. Classical $q$-Meixner orthogonality
then evaluates determinant specializations that are much smaller than the
uniform error in that limit. Combining these two scales gives the first
zero displacement for every fixed shift and the unequal displacements in
the second unshifted cluster. The finite-dimensional reduction is also an
exact identity: for fixed $r$, the size of the remaining determinant is
independent of $n$.

Several questions remain outside these results. The exact generic shifted
cyclotomic valuations are not determined by the guaranteed divisor.
Additional arithmetic content, different determinants, or sharper conjugate
estimates could change the norm comparison. For the analytic problem, the
individual splitting of clusters of multiplicity at least three requires
further information beyond the center value and first derivative.
Uniform zero displacement estimates when $r$ grows with $n$ likewise do
not follow from the fixed-$r$ determinant asymptotic.

The modular specializations in \cref{app:modular} are applications of
existing transcendence theorems. They are independent of the shift barrier
and the refined zero analysis, and are not claimed as newly resolved
irrationality problems.


\appendix
\section{Modular specializations: consequences of existing theory}\label{app:modular}

This appendix is logically independent of the determinant estimates. It
records applications of classical modular identities and transcendence
results, not a new proof of the underlying modular transcendence theorem.

For $p=e^{2\pi i\tau}$, $\operatorname{Im}\tau>0$, let
$J(p)=j(\tau)$. The theorem of Barr\'e-Sirieix, Diaz, Gramain, and
Philibert implies that $J(p)$ is transcendental for algebraic
$0<|p|<1$. The same conclusion follows from the stronger result in
\cite[Corollary 2]{Nesterenko}; that paper also identifies the earlier
1995 theorem in its historical discussion.
Weber's product and algebraic identity, in the normalization of
\cite[(9.9)--(9.11)]{Duke}, are
\begin{equation}\label{eq:weber}
 w(\tau)=\sqrt2\,p^{1/24}\prod_{n\ge1}(1+p^n),
 \qquad (w^{24}+16)^3=j(\tau)w^{24}.
\end{equation}
Only the twenty-fourth power is used, so the choice of root has no effect.

\begin{proposition}\label{prop:modular}
If $q\in\Qbar$ and $|q|>1$, then $E_q(1)$ is transcendental.
If $u\in\Qbar$, $|u|>1$, and $m\in\Z$, then
\begin{equation}\label{eq:special-family}
 E_{u^2}(\pm u^{2m+1})\quad\text{is transcendental}.
\end{equation}
\end{proposition}
\begin{proof}
Put $p=q^{-1}$, $\mu=E_q(1)$, and $z=p\mu^{24}$.
The nonvanishing product \eqref{eq:product} and \eqref{eq:weber} give
\begin{equation}\label{eq:j-product}
 J(p)=\frac{(1+256z)^3}{z}.
\end{equation}
If $\mu$ were algebraic, then so would be $z\ne0$ and $J(p)$,
contrary to the cited theorem.

To extend the family, splitting absolutely convergent products into even
and odd indices yields
\begin{equation}\label{eq:reciprocal}
 E_{u^2}(u)E_{-u}(1)=1,\qquad
 E_{u^2}(-u)E_u(1)=1.
\end{equation}
For example, writing $x=u^{-2}$, the first product is
\[
 \prod_{n\ge1}(1-x^{2n-1})\prod_{n\ge1}(1+x^n)=1,
\]
because $\prod(1+x^n)=\prod(1-x^{2n})/\prod(1-x^n)$.
The second identity follows by the same separation with $u^{-1}$.
Both reciprocal values are therefore transcendental.
Iterating \eqref{eq:functional} in both directions multiplies these
values by nonzero algebraic factors to obtain \eqref{eq:special-family}.
Indeed, each factor is $1\pm u^{2j+1}$ for some integer $j$, which is nonzero
because $|u|>1$.
\end{proof}

For rational $u=a/b>1$ in lowest terms, $q=u^2$ and
$\gamma_q=\log a/\log b$. Choosing $a=b+1$ makes this ratio approach one.
These are restricted pairs $(q,\alpha)$, not arbitrary independent rational
parameters. In particular, neither this appendix nor the shift theorem
asserts a new uniform critical value below $12/7$.

\section{Exact checks and reproducibility}\label{app:checks}

The proofs above are symbolic and analytic; finite computations are used
only to check signs, indices, and specializations. For the revised manuscript,
the following polynomial orders were checked by exact fraction-free
elimination after generating the entries from \eqref{eq:recurrence}.

\begin{table}[ht]
\centering
\caption{Exact orders at $q=0$ for $\alpha=2$ and $\mu=5$.}
\label{tab:checks}
\begin{tabular}{@{}rrrr@{}}
\toprule
$n$&$r$&Formula \eqref{eq:exact-order}&Computed order\\
\midrule
3&4&5&5\\
4&3&10&10\\
5&2&18&18\\
5&5&30&30\\
6&4&44&44\\
\bottomrule
\end{tabular}
\end{table}
\FloatBarrier

At $\alpha=2$, $\mu=3$, $n=5$, $r=2$, the order is $19$ instead of
the generic value $18$, as predicted by the vanishing lowest coefficient
in \cref{rem:exception}. The exact $\mu=0$ product was checked at
$q=2$, $\alpha=2$, $1\le n\le6$. The shifted dimension-reduction
identity was checked against direct rational Hankel determinants at
$q=2$, $1\le n\le4$, $0\le r\le4$.

These checks can be reproduced without additional data: form $v_0=\mu-1$,
apply the recurrence through index $2n-2+r$, and compute the indicated
determinant over $\Q[q]$ or $\Q$. To check the reduction identity, set
$\alpha=-2^{n+r+1}$ and $\mu=0$, compute its finite sums $S_N(a)$ in
rational arithmetic, and compare both sides after the stated normalization.
No numerical experiment is used to establish a transcendence assertion,
the location of an individual asymptotic zero, or the exactness of an
unspecified cyclotomic divisor. The single source file contains the entire
manuscript and bibliography and requires no external data files.

\begin{thebibliography}{9}
\raggedright
\bibitem{KRVZ}
C.~Krattenthaler, I.~Rochev, K.~V\"a\"an\"anen, and W.~Zudilin,
\emph{On the non-quadraticity of values of the $q$-exponential function and
related $q$-series}, Acta Arith. \textbf{136} (2009), no.~3, 243--269.
\href{https://doi.org/10.4064/aa136-3-4}{doi:10.4064/aa136-3-4};
\href{https://arxiv.org/abs/0812.2921}{arXiv:0812.2921}.

\bibitem{Rochev2011}
I.~P.~Rochev, \emph{On linear independence of values of certain $q$-series},
Izv. Math. \textbf{75} (2011), no.~1, 177--221.
\href{https://doi.org/10.1070/IM2011v075n01ABEH002531}
{doi:10.1070/IM2011v075n01ABEH002531}.

\bibitem{Macdonald}
I.~G.~Macdonald, \emph{Symmetric Functions and Hall Polynomials},
second edition, Oxford Mathematical Monographs, Clarendon Press, Oxford,
1995, Chapter I.
\href{https://doi.org/10.1093/oso/9780198534891.001.0001}
{doi:10.1093/oso/9780198534891.001.0001}.

\bibitem{KS}
R.~Koekoek and R.~F.~Swarttouw, \emph{The Askey-scheme of hypergeometric
orthogonal polynomials and its $q$-analogue}, Report 98-17,
Delft University of Technology, 1998, Section 3.13.
\href{https://tudelft.roelofkoekoek.nl/documents/as98.pdf}{1998 report};
\href{https://arxiv.org/abs/math/9602214}{arXiv:math/9602214}.

\bibitem{Nesterenko}
Yu.~V.~Nesterenko, \emph{Modular functions and transcendence questions},
Sb. Math. \textbf{187} (1996), no.~9, 1319--1348.
\href{https://doi.org/10.1070/SM1996v187n09ABEH000158}
{doi:10.1070/SM1996v187n09ABEH000158}.

\bibitem{Duke}
W.~Duke, \emph{Continued fractions and modular functions},
Bull. Amer. Math. Soc. (N.S.) \textbf{42} (2005), no.~2, 137--162.
\href{https://www.math.ucla.edu/~wdduke/preprints/bams4.pdf}{Author's preprint},
particularly (9.9)--(9.11).
\end{thebibliography}
\end{document}
